%%%%%%%%%%%%%%%%%%%%%%%%%%%%%%%%%%%%%%%%%%%%%%%%%%%%%%%%%%%%%%%%%%%%%%
% njuthesis 示例模板 v1.4.3 2025-05-21
% https://github.com/nju-lug/NJUThesis
%
% 贡献者
% Yu XIONG @atxy-blip   Yichen ZHAO @FengChendian
% Song GAO @myandeg     Chang MA @glatavento
% Yilun SUN @HermitSun  Yinfeng LIN @linyinfeng
% Yukai Chou @Muzimuzhi
%
% 许可证
% LaTeX Project Public License（版本 1.3c 或更高）
%%%%%%%%%%%%%%%%%%%%%%%%%%%%%%%%%%%%%%%%%%%%%%%%%%%%%%%%%%%%%%%%%%%%%%

%---------------------------------------------------------------------
% 一些提升使用体验的小技巧：
%   1. 请务必使用 UTF-8 编码编写和保存本文档
%   2. 请务必使用 XeLaTeX 或 LuaLaTeX 引擎进行编译
%   3. 不保证接口稳定，写作前一定要留意版本号
%   4. 以百分号(%)开头的内容为注释，可以随意删除
%---------------------------------------------------------------------

%---------------------------------------------------------------------
% 请先阅读使用手册：
% http://mirrors.ctan.org/macros/unicodetex/latex/njuthesis/njuthesis.pdf
%---------------------------------------------------------------------

\documentclass[
    % 模板选项(注意右端逗号)：
    %
    type = master, % 文档类型，默认为本科生
    degree = academic,        % 学位类型，默认为学术型
    %
    % nl-cover,   % 是否需要国家图书馆封面，默认关闭
    % decl-page,  % 是否需要诚信承诺书或原创性声明，默认关闭
    %
    %   页面模式，详见手册说明
    % draft,                  % 开启草稿模式
    % anonymous,              % 开启盲审模式
    % minimal,                % 开启最小化模式
    %
    %   单双面模式，默认为适合印刷的双面模式
    % oneside,                % 单面模式，无空白页
    % twoside,                % 双面模式，每一章从奇数页开始
    %
    %   字体设置，不填写则自动调用系统预装字体，详见手册
    fontset = fandol % |mac|macoffice|win|none,
  ]{njuthesis}

% 模板选项设置，包括个人信息、外观样式等
% 较为冗长且一般不需要反复修改，我们把它放在单独的文件里
% njuthesis 参数设置文件 v1.4.3 2025-05-21

% 一些提醒：
%   1. \njusetup 内部千万不要有空行
%   2. 使用英文半角逗号(,)分隔选项
%   3. 等于号(=)两侧的空格会被忽略
%       3.1. 为避免歧义，请用花括号({})包裹内容
%   4. 本科生无需填写的项目已被特别标注
%   5. 可以尽情删除本注释

% info 类用于录入个人信息
%   带*号的为对应英文字段
\njusetup[info]{
    title = {信息激励与情感驱动：捐赠信息披露和社区讨论情感对比特币开源开发者活动的影响机制},
    % 中文题目
    % 直接填写标题就是自动换行
    % 可以使用换行控制符(\\)手动指定换行位置
    %
    title* = {Information Incentives and Emotional Drives: The Impact Mechanism of Donation Information Disclosure and Community Discussion Sentiment on Bitcoin Open-Source Developer Activity},
    % 英文题目
    %
    author = {冉红云},
    % 作者姓名
    %
    author* = {HongyunRan},
    % 作者英文姓名
    % 一般使用拼音
    %
    keywords = {比特币,开源社区,信息披露,开发者激励,开发者情感},
    % 中文关键词列表
    % 使用英文半角逗号(,)分隔
    %
    keywords* = {Bitcoin,{Open-Source Community},{Information Disclosure},{Developer Motivation},{Developer Sentiment}},
    % 英文关键词
    % 使用英文半角逗号(,)分隔
    %
    grade = {2023},
    % 年级
    %
    student-id = {522023140076},
    % 学号或工号
    % 研究生请斟酌大小写字母格式
    % 本模板并不会自动更正大小写
    %
    department = {信息管理学院},
    department* = {School of Information Management},
    % 院系
    %
    major = {图书情报},
    major* = {Library and Information Studies},
    % 专业
    %
    % major = {封面专业,摘要专业},
    % 研究生专业型学位可能遇到两处内容不一致的情况
    %
    supervisor = {颜嘉麒,教授},
    supervisor*= {JaqiYan,Professor},
    % 导师全称
    % 使用英文半角逗号(,)分隔中文姓名和职称
    %
    % supervisor-ii = {第二导师姓名,副教授},
    % supervisor-ii* = {Associate Professor My Second Supervisor},
    % 第二导师全称
    % 如果确实没有第二导师，不填写即可
    %
    submit-date = {2026-05-19},
    % 提交日期
    % 格式为 yyyy-mm-dd
    % 不填就是编译当天日期
    %
    %
    % 以下均为研究生项
    %
    % degree = {工程硕士},
    % degree* = {Master of Engineering},
    % 覆盖默认学位名称
    %
    field = {开源社区治理},
    field* = {Open source community developer incentives},
    % 研究领域
    %
    chairman = {康乐乐~教授},
    % 答辩委员会主席
    % 推荐使用波浪号(~)分隔姓名和职称
    %
    reviewer = {
        柯青~教授,
        赵月华~副教授
    },
    %
    % 答辩委员会成员
    % 一般为四名，使用英文半角逗号(,)分隔
    %
    clc = {G350},
    % 中国图书分类号
    %
    udc = {001.18},
    % 国际图书分类号
    %
    secret-level = {},
    % 密级
    %
    defend-date = {2026-05-24},
    % 答辩日期
    % 格式为 yyyy-mm-dd
    % 不填就是编译当天日期
    %
    email = {522023140076@smail.nju.edu.cn},
    % 电子邮箱地址
    % 只用于出版授权书
    %
    %
    % 以下用于国家图书馆封面
    confer-date = {2026-06-01},
    % 学位授予日期
    %
    bottom-date = {2026-06-01},
    % 封面底部日期
    %
    supervisor-contact = {
        南京大学~
        江苏省南京市栖霞区仙林大道163号
    },
    % 导师联系方式
}

% bib 类用于参考文献设置
\njusetup[bib]{
    resource = {thesis.bib},
    option = {
        sorting = gb7714-2015,
        gblanorder = englishahead,
    }
}

% image 类用于载入外置的图片
\njusetup[image]{
    % path = {{./figure/}{./image/}},
    % 图片搜索路径
    %
    % nju-emblem = {nju-emblem},
    % nju-name = {nju-name},
    % 校徽和校名图片路径
    % 建议使用 PDF 格式的矢量图
    % 使用外置图片有助于减少编译时间
    % 空置时会自动使用 njuvisual 宏包绘制
    %
    nju-emblem = {nju-emblem-purple},
    nju-name = {nju-name-purple},
    % 替换为紫色版本
    % 这个选项只能填写一次
    % 切换时要注释掉上方的黑色版本
}

% abstract 类用于设置摘要样式
\njusetup[abstract]{
    toc-entry = false,
    % 摘要是否显示在目录条目中
    %
    % underline = false,
    % 研究生英文摘要页条目内容是否添加下划线
    %
    % title-style = strict|centered|natural
    % 研究生摘要标题样式，详见手册
}

% 目录自身是否显示在目录条目中
\njusetup{
    % tableofcontents/toc-entry = false,
    % listoffigures/toc-entry   = false,
    % listoftables/toc-entry    = false
}

% 为目录中的章标题添加引导线
\njusetup[tableofcontents/dotline]{chapter}

% math 类用于设置数学符号样式，功能详见手册
\njusetup[math]{
    % style              = TeX|ISO|GB,
    % 整体风格，缺省值为国标（GB）
    % 相当于自动设置以下若干项
    %
    % integral           = upright|slanted,
    % integral-limits    = true|false,
    % less-than-or-equal = slanted|horizontal,
    % math-ellipsis      = centered|lower,
    % partial            = upright|italic,
    % real-part          = roman|fraktur,
    % vector             = boldfont|arrow,
    % uppercase-greek    = upright|italic
}

% theorem 类用于设置定理类环境样式，功能详见手册
\njusetup[theorem]{
    % define,
    % 默认创建内置的七种定理环境
    %
    % style         = remark,
    % header-font   = \sffamily \bfseries,
    % body-font     = \normalfont,
    % qed-symbol    = \ensuremath { \male },
    % counter       = section,
    % share-counter = true,
    % type          = {...},
    % define,
    % 以上设置项在重新调用 theorem/define 后生效
}

% footnote 类用于设置脚注样式，功能详见手册
\njusetup[footnote]{
  style = circled,
  % 使用圈码编号
  hang = false
}

% 页眉页脚内容设置
\njusetup{
  header/content = {
      {OR}{\rightmark},   % 奇数页页眉右侧：章名
      {OL}{},             % 奇数页页眉左侧：空
      {EL}{\leftmark},    % 偶数页页眉左侧：节名
      {ER}{}              % 偶数页页眉右侧：空
    },
  footer/content = {
      {C}{\thepage}       % 页脚居中：页码
    }
}

% 页眉页脚的字体样式
\njusetformat{header}{\small\kaishu}
\njusetformat{footer}{}

% 在盲审模式下隐藏学校信息
% \njusetup{anonymous-mode/no-nju}

% 一些灵活调整
% \njusetname{type}{本科毕业设计}                 % 我做的是毕业设计
% \njusetname{notation}{术语表}                   % 更改符号表名称
% \njusetlength{crulewd}{240pt}                   % 加长封面页下划线
% \njusetformat{subsection}{\normalfont\bfseries} % 修改 subsection 为小四号粗体宋体
% \njusetformat{tabular}{\zihao{-4}\bfseries}     % 修改表格环境的字号
% \EditInstance{nju}{u/cover/emblem-img}{align=l} % 左对齐的本科生封面校徽

% 根据用户要求添加的格式设置
% 1. 摘要格式设置
\njusetup[abstract]{
  title-style = centered,  % "摘要"二字居中
}

% 2. 正文标题格式设置
\njusetformat{chapter}{\centering\heiti\zihao{-2}\bfseries}  % 一级标题：黑体，小二，加粗，居中
\njusetformat{section}{\heiti\zihao{-3}\bfseries}            % 二级标题：黑体，小三，加粗
\njusetformat{subsection}{\heiti\zihao{4}\bfseries}          % 三级标题：黑体，四号，加粗

% 3. 正文格式设置
\renewcommand{\baselinestretch}{1.5}  % 1.5倍行距
\zihao{-4}\songti                     % 正文：小四号宋体

% 4. 脚注格式设置
\njusetup[footnote]{
  style = circled,
  hang = false
}

% 5. 参考文献格式设置（在biblatex中通过其他方式设置）
% \renewcommand{\bibfont}{\zihao{5}\songti}  % 五号字体，中文宋体，英文Times New Roman

% 6. 关键词格式设置（在abstract环境中自动处理）
% 7. 页码格式已在njuthesis.cls中设置（摘要罗马数字，正文阿拉伯数字）


% 解决 hyperref PDF 书签中的字体命令警告
\pdfstringdefDisableCommands{%
  \let\heiti\relax
  \let\songti\relax
  \let\kaishu\relax
  \let\zihao\relax
  \let\bfseries\relax
  \let\centering\relax
}

% 允许更宽松的断行，减少 overfull/underfull 警告
\sloppy

% 自行载入所需宏包
\usepackage{booktabs}    % 三线表（toprule/midrule/bottomrule）
\usepackage{tabularx}    % 自适应列宽表格
\usepackage{multirow}    % 跨行合并单元格
% \usepackage{subcaption} % 嵌套小幅图像，比 subfig 和 subfigure 更新更好
% \usepackage{siunitx} % 标准单位符号
% \usepackage{algorithm,algorithmic} % 展示算法伪代码

% 在导言区随意定制所需命令
\NewDocumentCommand\todo{m}{\textbf{[待补充：#1]}} % 待填写标记

\begin{document}

%---------------------------------------------------------------------
%   封面、摘要、目录
%---------------------------------------------------------------------

\maketitle
\raggedbottom

% 中文摘要
\begin{abstract}
比特币是目前市值最高的加密货币，其开发完全依赖全球开源开发者的自愿贡献。
如何激励开发者持续贡献，是比特币社区及众多开源项目面临的核心挑战。
既有研究主要关注实际货币激励对开发者行为的影响，而对于
\textbf{捐赠信息披露}本身所产生的激励效应，以及
\textbf{社区内部情感状态}对开发者活动的动态影响，尚缺乏系统研究。

本文以 GitHub 上的比特币开源项目为研究对象，聚焦两个相互补充的研究问题：

其一，外部非盈利基金会（Bitcoin Brink）的捐赠信息披露如何影响开发者活动。
Brink 在运营过程中存在两次信息披露：在 X 平台公布收到的捐款信息（非定向披露），
以及在官网公布受资助开发者名单（定向披露）。
本文采用面板数据固定效应模型与工具变量方法，
在剔除实际受资助开发者后，识别信息披露的纯信息效应。
研究发现：两种披露均显著正向影响开发者贡献，
但定向披露对非定向披露的正向效应产生负向调节，
印证了期望失验理论所预测的失望效应。

其二，比特币社区 Issue 讨论中的情感状态如何影响开发者活动，
以及代码活动与讨论活动之间是否存在替代效应。
本文采用向量自回归（VAR）模型，发现滞后一期的负面情感正向影响代码活动与讨论活动，
滞后两期的代码活动负向影响讨论活动，
揭示了"情感→代码/讨论→讨论"的影响路径。
上述结果分别以情感即信息理论和目标手段替代理论加以解释。

本研究从信息视角重新审视开源开发者的激励机制，
为比特币及其他开源社区的捐赠信息披露策略设计提供实证参考，
并对图书情报学中的信息行为与用户激励研究作出理论贡献。
\end{abstract}

% 英文摘要
\begin{abstract*}
Bitcoin, the world's highest-valued cryptocurrency, relies entirely on voluntary contributions
from open-source developers worldwide. Sustaining developer motivation remains a central
challenge for the Bitcoin community and open-source projects at large. While prior research
has predominantly examined the direct effects of monetary incentives, the motivating impact
of \textit{donation information disclosure} itself—decoupled from actual financial transfers—
and the role of \textit{community sentiment} in shaping developer behavior remain underexplored.

This thesis examines two complementary research questions using the Bitcoin GitHub repository
as the empirical context. First, it investigates how two distinct donation disclosure
strategies employed by the non-profit organization Bitcoin Brink affect developer activity:
untargeted announcements of received donations posted on platform X, and targeted
grantee disclosures published on Brink's official website. Employing panel fixed-effects
models and instrumental variable estimation—while excluding actual grantees to isolate
pure information effects—the study finds that both disclosure types positively influence
developer contributions. However, targeted disclosure negatively moderates the positive
effect of untargeted disclosure, consistent with the disappointment mechanism predicted
by expectancy disconfirmation theory.

Second, the thesis examines how community sentiment expressed in Issue discussions
dynamically influences developer activities, and whether substitution effects exist between
code and discussion activities. Using Vector Autoregression (VAR) models, the analysis
reveals that lagged negative sentiment positively influences both code and discussion
activities, while increased code activity subsequently suppresses discussion activity.
This yields the influence pathway: sentiment → code/discussion → discussion, explained
through affect-as-information theory and goal–means substitution theory.

This research reframes open-source developer motivation from an informational perspective,
offering actionable insights for donation disclosure strategy design in Bitcoin and
broader open-source communities, and contributing theoretically to information behavior
research in library and information science.
\end{abstract*}

% 目录
\tableofcontents

%---------------------------------------------------------------------
%   正文
%---------------------------------------------------------------------
\mainmatter

%%====================================================================
\chapter{绪论}
%%====================================================================

\section{研究背景}

过去数十年间，开源软件（Open-Source Software，OSS）已从小众的技术运动成长为
现代数字基础设施不可或缺的组成部分\cite{lerner2002simple}。
Linux 操作系统支撑着全球绝大多数服务器与移动设备，
Apache HTTP Server 托管着互联网上数量庞大的网站，
Python、Git 等工具深度嵌入全球数百万开发者的日常工作流程——
这些广为人知的软件均是开源模式下协作开发的产物\cite{lakhani2004open}。
与商业软件的封闭开发不同，开源项目以源代码公开、全球开放协作为基本特征，
任何人均可查看代码、提交改进、报告问题，无需与项目组织存在任何雇佣关系。
这种模式极大地降低了知识传播的门槛，形成了高度去中心化的创新生态；
与此同时，也带来了一个根本性的挑战：既然贡献是自愿的，
那么究竟是什么驱使着全球各地的开发者持续投入时间和精力？

这一问题在比特币（Bitcoin）社区中表现得尤为突出。
比特币是 Satoshi Nakamoto 于2008年提出的点对点电子现金系统\cite{nakamoto2008bitcoin}，
其底层协议自诞生之日起便以完全开源的形式运作，
项目代码库托管于 GitHub 平台，目前已积累超过76,000名关注者，
吸引了来自全球的数千名开发者贡献代码与文档。
比特币不同于其他开源项目的关键在于：它背后没有任何公司主体为开发活动提供资金保障，
也没有雇主为核心开发者发放薪酬。整个项目的存续与演进完全建立在志愿贡献的基础之上，
而这些贡献直接关系到一个数千亿美元规模的金融网络的安全性与可靠性。
一旦核心开发者离去或社区贡献萎缩，比特币网络将面临严峻的技术维护风险。
因此，如何有效激励开发者持续贡献，是比特币社区乃至整个开源生态面临的核心治理议题。

学术界对开源开发者激励机制的研究已积累了丰富成果。
早期研究发现，开发者的参与主要由内在动机驱动：对技术问题的纯粹兴趣\cite{shah2006motivation}、
在同行社区中积累声誉的渴望\cite{lerner2002simple}，以及通过贡献提升个人技能、
从而在劳动力市场中获得更好机会的职业考量\cite{fang2009understanding}，
共同构成了驱动开发者持续投入的基础动力。
此外，Davidson 等人\cite{davidson2014older}指出，利他精神与自我实现感同样是重要的内在激励来源。
这些研究共同刻画了一个图景：在开源生态中，
开发者的自主性与社区归属感往往是持续贡献的核心支撑，
而非来自外部的物质回报。

然而，随着开源软件在商业和社会中的重要性持续上升，
外部经济激励机制也开始进入研究者的视野。
货币奖励被证明能够在一定程度上促进开发者的贡献\cite{wang2022monetary}，
但其与内在动机的关系颇为微妙——若管理不当，不适当的金钱激励反而可能产生"挤出效应"，
削弱开发者原本的内在驱动力\cite{benabou2003intrinsic,qiao2021mitigating}。
2019年，GitHub 正式推出 Sponsors 机制，首次为开源贡献提供了平台级的捐赠基础设施，
使任何人都可以向自己欣赏的开发者提供经济支持。
这一机制的出现引发了学界的广泛关注。
Shimada 等人\cite{shimada2022github}的研究发现，
部分开发者将赞助本身视为一种新型的社区贡献方式；
Zhang 等人\cite{zhang2022who}进一步揭示，赞助在短期内对开发活动具有温和的正向效应，
但获得赞助的可能性与开发者在社区中的社会地位密切相关；
Conti 等人\cite{conti2023incentivizing}则通过更长期的追踪研究发现，
实际收到赞助反而可能对创新活动产生长期负向影响，
这与挤出效应的理论预测相吻合。
Overney 等人\cite{overney2020not}的研究也表明，
单纯依靠捐赠机制提升开源贡献的效果并不显著，
暗示捐赠的效用可能远比表面看起来复杂。

上述研究的密集出现，折射出开源生态在现实层面所面临的深层可持续性困境。
现代数字经济的高速运转建立在大量关键开源组件之上，然而这些组件的维护工作往往
高度集中于少数志愿开发者，缺乏稳定的制度性支持。
2021年底曝光的 Apache Log4j 漏洞事件以极为惨烈的方式揭示了这一隐忧——
一个嵌入全球数十亿软件系统的核心日志库，长期仅由数名志愿者以业余时间维护，
其安全漏洞的存在几乎使整个互联网基础设施陷入系统性风险\cite{petryk2023gharchive}。
在比特币生态中，这种脆弱性以更为独特的方式呈现：
比特币网络目前承载着超过万亿美元规模的资产，
而其底层协议的核心维护工作却高度依赖为数不多的全职开发者，
任何核心人才的流失或社区参与的持续萎缩，
都有可能对协议演进能力乃至网络整体安全产生深远影响。
正是在这一背景下，以 Brink 为代表的专项资助机构应运而生，
探索出基于信息透明化和分阶段披露的独特激励路径，
以期在不违背去中心化治理理念的前提下，维持更广泛开发者群体的参与意愿与贡献活力。
对这一机制的科学评估，因此具有超越单一案例的现实普遍意义。

综合上述研究可以发现，现有文献对捐赠的理解主要集中于其作为货币转移的经济效应，
即"钱到位了，开发者是否更积极"这一层面。
然而，在实际的捐赠生态中，信息的流动往往先于资金的转移，
而且覆盖的受众也远比实际受益者更广——
当一个资助机构在社交媒体上宣布收到了一笔捐款时，
所有关注该机构的开发者都会接收到这条信息，
无论他们最终是否能从中受益。这条信息本身会不会影响他们的行为？
这种仅凭信息而非真金白银的激励效应，
在既有文献中几乎没有得到系统性的研究与检验。
这一研究空白尤为值得关注，因为信息披露相较于实际资金分配具有近乎为零的边际成本，
同时又能触达整个社区而非少数受资助者，
若其激励效应确实存在，则将为开源社区的低成本治理提供重要的政策工具。
从信息经济学的视角审视，捐赠公告的本质是一种信号传递行为，
它向潜在贡献者同时传递了项目获得外部认可、机构具有奖励意愿以及未来奖励可得性等多重信号，
这些信号会在开发者心中激活或调整预期结构，进而影响其参与决策\cite{huang2021examining}。
类似的逻辑在众筹研究中已有部分印证——项目进展公告能够显著提升后续的支持率，
即便公告内容并不伴随任何即时的经济转移\cite{huang2021examining}。
然而，在开源开发者群体中，信息披露能否以类似机制触发可测量的行为变化，
尤其是不同类型披露（普惠型与定向型）是否产生差异化效应，
目前在实证层面仍是空白。

与此同时，在社区内部的信息维度上，同样存在一个亟待填补的研究空白。
开源项目的 Issue 讨论区不仅是技术问题协商的场所，
也是社区成员情感状态汇聚与传导的媒介。
当项目遭遇重大 Bug、核心功能争议或对外部环境变化的集体焦虑时，
Issue 区的讨论情感会随之出现明显的极性变化。
Guzzi 等人\cite{guzzi2013}发现，开源社区的协作沟通中蕴含着丰富的情感信息，
这些信息构成了社区健康状态的重要信号。
然而，社区情感如何在时间维度上影响开发者的具体行为，
情感变化是否能够预测随后的代码活动或讨论活动，
以及不同类型的开发活动之间是否存在相互促进或相互抑制的动态关系，
目前尚缺乏系统性的实证探讨。
现有将情感分析引入开源社区的研究，大多以描述情感分布的静态特征
或考察情感与代码质量的截面相关性为主要目标\cite{guzzi2013}，
较少采用动态因果推断框架，系统地检验情感冲击在时间轴上如何传导至开发者的行动决策。
更重要的是，代码活动与讨论活动之间的动态互动——
当开发者通过编写代码回应社区问题时，是否会相应地减少讨论参与，
形成一种活动层面的替代关系——这一机制既有理论依据，
又对理解开源社区的协作动态具有重要意义，
却在实证研究中几乎未曾得到正面检验。
在信息行为研究的视角下，情感本身就是一种重要的信息资源，
它向社区成员传递项目状态的信号，影响其行动意愿与行动方向。
从这一视角出发理解开发者行为，有可能揭示出超越个体激励机制的社区层面动态规律。

正是在上述背景下，本研究将视线转向比特币开源社区，
尝试从信息视角对上述两个研究空白进行系统性的实证回应。

\section{研究场景：比特币社区与 Brink 基金会}

比特币开源社区为本研究提供了几乎难以在其他情境中复现的理想研究场景，
核心原因在于一个独特的研究对象：非盈利组织 Bitcoin Brink（以下简称"Brink"）
及其所形成的两阶段捐赠信息披露机制。

Brink 成立于2020年，是目前规模最大的专注于比特币核心协议开发的资助机构，
采用众筹模式面向个人捐赠者和机构投资者募集资金，
并通过严格的申请审核流程为符合条件的比特币开发者提供全职工作资助。
由于比特币本身的去中心化哲学，Brink 刻意保持独立的非盈利定位，
其资金来源和资助决策均对外公开，以维护社区信任。
正是这种对透明度的高度重视，催生了一套具有内在研究价值的两阶段信息披露机制。

第一阶段，当 Brink 收到捐款时，会在其 X 平台（原 Twitter）官方账号上发布公告，
公开感谢捐赠方并注明捐款金额，但不说明这笔资金将用于资助哪位具体的开发者。
这类公告面向所有关注者广泛传播，每一位比特币开发者都可能看到这条信息；
由于受益人尚未确定，公告中所传递的核心信号是"这个社区正在获得外部支持，
有资金可以用于奖励开发者"。本研究将这类披露定义为\textbf{非定向披露（Untargeted Disclosure）}，
强调其普惠性——每位开发者均可合理地将自己视为潜在受益者。
第二阶段，经过通常为数周乃至数月的申请与审核流程后，
Brink 会在其官方网站上正式宣布当期受资助开发者名单，
列出姓名、资助金额及资助期限，并对其贡献做出肯定性说明。
与第一阶段不同，这类披露所传递的核心信号是"谁因为贡献突出而获得了认可与资金支持"。
本研究将其定义为\textbf{定向披露（Targeted Disclosure）}，
因为它指向具体的个人，对于未入选者而言也意味着明确的排除。

两次披露在时间上相互分离、在内容上存在本质差异，
这在方法论层面具有重要价值：研究者得以将两种信息效应分别识别，
而无需将其与实际资金的到达混为一谈。
更重要的是，在本研究的观测期内（2022年1月至2023年12月），
Brink 实际资助的开发者数量极为有限——两年内仅资助了4名开发者转为全职从事比特币开发。
这意味着，通过将这4名实际受资助者从分析样本中剔除，
研究者可以专注于研究对于其余1613名开发者而言，
捐赠信息的公开本身——而非任何真实的金钱转移——
是否对其贡献行为产生了可测量的影响。
这种"信息与金钱的干净分离"，在开源捐赠研究中极为罕见，
构成了本研究得以回答核心问题的关键条件。

除了捐赠信息披露机制之外，比特币社区还具备另一项研究优势：
完整、开放且可系统处理的 Issue 讨论历史。
比特币项目自创立以来积累了大量的 Issue 记录，
涵盖 Bug 报告、功能讨论、安全议题乃至社区内部的技术争论，
这些文本构成了追踪社区情感动态变化的完整语料库。
结合 GH Archive 对 GitHub 活动数据的完整记录、X 平台的公告历史、
Brink 官网的资助档案以及 Google Trends 的搜索量数据，
本研究得以在同一社区中、围绕同一组开发者群体，
同时探讨外部信息与内部情感两个维度对开发者行为的影响。
数据来源的多元性与可靠性，是本研究最终选择比特币社区作为研究对象的重要原因。

\section{研究问题}

建立在上述背景与研究场景之上，本研究聚焦两个相互补充的核心问题。

第一个研究问题关注外部信息对开发者行为的影响。
Brink 基金会在实践中形成了非定向披露与定向披露两种信息传播方式，
二者在受众结构、信号内容和心理机制上存在本质差异，
却可能产生相互交织的激励效果。本研究希望厘清：

\textbf{RQ1}：Bitcoin Brink 基金会的非定向捐赠披露与定向捐赠披露如何分别影响
开发者的代码贡献、审查贡献和讨论贡献？当两种披露同时存在时，
二者之间是否产生交互效应？

第二个研究问题则将目光转向社区内部，探讨情感信息的传导机制。
社区讨论中的情感状态不仅是开发者体验的反映，
也可能作为一种项目状态的信号，影响开发者的行动选择。
在代码贡献与讨论参与这两种典型的开发活动之间，是否存在情感所触发的动态路径？

\textbf{RQ2}：比特币开源社区 Issue 讨论中的情感状态如何在时间维度上影响
开发者的代码活动与讨论活动？代码活动的变化是否进一步影响讨论活动，
形成可识别的跨活动传导路径？

\section{研究方法概述}

针对两个研究问题，本研究采取差异化的计量策略。

研究一所面对的核心方法论挑战是内生性问题：
开发者的活动水平可能反向影响外部捐款规模，
进而影响 Brink 的披露行为，从而在观测数据中产生反向因果偏误。
本研究采用\textbf{面板数据固定效应模型（Panel Fixed-Effects Model）}
作为基准识别框架，通过控制个体固定效应与时间固定效应，
消除不可观测的个体异质性与共同时间趋势的干扰。
在此基础上，进一步引入\textbf{工具变量方法（Instrumental Variable，IV）}，
以"Brink Bitcoin"关键词在 Google Trends 上的搜索量指数
作为非定向披露金额的工具变量，以外生的公众关注度变化来识别披露金额的因果效应，
解决反向因果导致的估计偏差。

研究二面对的核心方法论挑战是变量间的双向因果关系：
情感、代码活动和讨论活动三者之间可能存在相互影响，
无法在单方程框架下合理处理。
本研究采用\textbf{向量自回归模型（Vector Autoregression，VAR）}，
将三个变量均作为内生时间序列联合建模，
通过格兰杰因果检验（Granger Causality Test）在统计层面识别变量间的预测关系，
并借助脉冲响应函数（Impulse Response Function，IRF）
直观呈现某一变量受到冲击后其他变量的动态响应路径，
从而系统描绘社区情感与开发活动之间的动态传导机制。

\section{研究意义}

\subsection{理论意义}

本研究最核心的理论贡献，在于从信息视角重新审视开源开发者的激励机制，
将既往研究中混同于货币效应的信息效应单独提炼出来加以检验。
在已有的开源激励研究中，捐赠通常被视为一种货币转移行为，
其影响机制被归结为对开发者的直接经济补偿\cite{zhang2022who,wang2022influence}。
然而，在实际运作中，捐赠信息的公开传播往往先于、且波及范围远大于真实的资金分配，
形成了一种超越货币本身的信号效应。
本研究通过对样本的精心处理与工具变量的设计，
将这一信号效应从货币效应中剥离，首次在开源社区情境下提供了
"捐赠信息披露本身是否构成有效激励因素"的实证证据，
从而丰富了信息激励理论的层次与应用边界。

在理论框架层面，本研究创新性地将\textbf{期望理论（Expectancy Theory）}
与\textbf{期望失验理论（Expectancy Disconfirmation Theory）}联合应用于
同一研究情境的不同阶段。非定向披露通过激活"潜在受益"期望影响开发者动机，
而定向披露的公布随即在大多数非受助者中引发了期望落空，
两种机制在时间上先后作用，共同塑造了开发者对信息披露的复合响应。
这种基于期望认知动态的分析框架，对于理解其他需要区分"普遍性激励"与
"选择性激励"效果的管理情境同样具有理论借鉴价值。

与此同时，对于社区情感影响路径的研究，将\textbf{情感即信息理论（Affect-as-Information Theory）}
引入开源开发者行为研究领域。
情感作为信息的视角强调，情感状态本身携带着关于当前处境的认知内容，
引导个体调整判断与行动方向\cite{schwarz1983moods}。
将这一框架应用于开源社区，意味着把 Issue 讨论中的集体情感
理解为一种社区状态的公共信号，而非单纯的情绪宣泄。
此外，\textbf{目标手段替代理论（Goal-Means Substitution Theory）}被用于解释
代码活动对讨论活动的抑制效应，这一跨理论的应用拓展了两种理论在数字协作情境中的适用范围，
并为理解开源社区中不同类型活动之间的内在关联提供了新的理论工具。

从图书情报学的学科视角出发，本研究将"信息行为"这一核心议题延伸至
开源软件社区这一新兴的数字协作场景，
将开发者对信息披露的响应与情感信号的加工理解为一种特殊的信息利用与决策行为。
在方法论层面，本研究融合了面板计量经济学与时间序列分析方法，
为图书情报学的实证研究提供了多源异构数据综合分析的示范案例，
有助于推动本学科在量化研究方法上的多元化发展。

在研究设计层面，本研究构建了"信息效应与货币效应分离"的识别框架，
通过将实际受资助者从样本中剔除、并以工具变量方法处理内生性问题，
在开源社区研究情境中实现了较为严格的因果推断。
这种将准自然实验思想系统融入开源行为研究的方法论路径，
有别于现有文献中普遍采用的相关性分析范式。
对于希望从"什么因素与开发者行为相关"推进到
"什么因素真正驱动了开发者行为"的后续研究者而言，
本研究提供了一个可供参照的方法论示范，
有助于推动开源社区研究从描述性研究向因果推断研究的范式演进。

\subsection{实践意义}

对于以 Brink 为代表的开源基金会而言，本研究揭示了信息披露方式的设计对于
社区开发活力的实际影响，这一发现具有直接的政策参考价值。
非定向披露能够在整个开发者群体中激活期望，形成广泛的动机激励，
适合定期发布以维护社区的整体参与感；
但定向披露在明确受益人之后，会不可避免地在未入选者中产生失落效应，
从而削弱非定向披露的长期激励效果。
这提示基金会在进行信息披露时，应当有意识地管理披露的节奏与表述方式，
例如在定向披露公告中着重强调资助计划的持续性和未来的更多机会，
以缓解未入选者的期望落空感。

对于比特币社区乃至更广泛的开源生态而言，
社区情感对开发活动的预测作用，为社区健康管理提供了一个实践启示：
情感极性的下滑可以作为社区面临冲突或技术困境的预警信号，
在这一时期，代码活动往往会随之上升，说明社区成员正在自发地调动资源解决问题，
社区管理者可以借助情感监测系统识别此类时间节点，
及时提供协调资源、降低开发者面临的障碍。

更广泛地看，本研究所揭示的"以信息换激励"逻辑，
对于所有依赖自愿贡献、且资源有限的数字公共品社区均具有参考意义。
在无法依靠大规模直接资助的现实约束下，
信息的透明化与策略性披露可以作为货币激励的低成本替代或有益补充，
在维持社区信任与凝聚力的同时，对更广泛的潜在贡献者产生持续的参与激励。
这一视角不仅适用于加密货币生态，
对于同样面临可持续性挑战的维基百科类志愿知识贡献平台、
公民科学社区以及政府推动的开放数据生态，同样具有一定的政策启发价值。

\section{论文结构安排}

本文共七章，各章内容安排如下：第二章系统梳理相关理论基础与文献综述，
依次涵盖开源开发者激励机制的研究脉络、捐赠与赞助机制的文献进展、
信息披露的理论基础，以及期望理论、期望失验理论、情感即信息理论等核心理论框架的回顾。
第三章作为两项实证研究的共同基础，统一呈现整体研究框架与设计逻辑、
六类数据源的采集路径与融合流程、两种方法论的选择依据，
以及研究过程中的数据伦理考量。
第四章为研究一，围绕捐赠信息披露对开发者活动的影响，
完整呈现研究场景与假设推导、数据来源与变量定义、计量模型构建，
以及主效应结果、稳健性检验和机制分析的实证发现。
第五章为研究二，聚焦比特币社区情感对开发者活动的动态影响，
介绍 VAR 模型的设定与估计过程，汇报格兰杰因果检验、
脉冲响应分析与方差分解的结果，并对影响路径作出理论解释。
第六章在两项独立研究的基础上进行综合讨论，提炼理论贡献与实践启示，
并客观说明本研究的局限性及未来可改进的方向。
第七章总结全文的主要研究结论，并对后续研究的拓展方向提出展望。


%%====================================================================
\chapter{理论基础与文献综述}
%%====================================================================

\section{开源软件开发与开发者激励}

\subsection{开源软件的特征与发展历程}

开源软件（Open-Source Software, OSS）是指以源代码公开为基本特征、
允许任何人自由查看、使用、修改和分发的软件\cite{lakhani2004open}。
与商业软件的封闭开发模式不同，开源项目的开发过程本质上是一种以互联网为媒介的
分布式协作——参与者来自全球不同机构，通常不受雇佣合同约束，
以自愿形式贡献代码、文档和讨论，彼此之间的协调依靠版本控制系统、
邮件列表和代码审查流程等技术与社会机制来实现。

自由软件运动在20世纪80年代由 Richard Stallman 奠定理念基础，
而 Linux 内核于1991年的出现则将这一理念推向工程实践的前沿。
进入21世纪，互联网基础设施的快速普及与 Git 版本控制工具的成熟，
极大降低了全球开发者协作的组织成本，以 GitHub 为核心的开源生态随之形成。
时至今日，Linux 内核支撑着全球绝大多数服务器与移动设备，
Python 成为数据科学和人工智能研究的通用语言，
这些产品背后均是数以千计的志愿开发者多年协作的产物。

从经济学角度审视，开源软件具有公共品的典型特征：
非竞争性与非排他性共同导致潜在的"搭便车"问题——
在理论上，人人都倾向于坐享他人贡献而不愿自行付出成本\cite{lerner2002simple}。
然而，开源软件生态的实际状况与这一悲观预测相悖：
大量高质量项目维持了数十年的持续演进，吸引了遍布全球的贡献者。
这一"反常"现象激发了学界对开源开发者动机机制的持续探索，
形成了一套有别于传统经济人假设的理论体系。

值得关注的是，开源软件的经济重要性在近十余年间经历了质的跃升。
据估计，企业在其核心软件产品中使用的代码超过70\%来源于开源组件，
全球开源软件每年为商业应用贡献的价值已达数千亿美元规模。
与此同时，维持这一巨大价值的开发者群体却在规模上极度有限，
且整体上处于志愿贡献的状态——许多关键项目的核心维护工作仅集中在个位数的开发者身上。
这种"巨大经济价值与极少志愿维护者"之间的不对称，
使得理解和维持开发者贡献动机成为一个兼具学术价值和现实紧迫性的重要议题。
近年来，学界开始广泛使用"开源可持续性（open source sustainability）"
这一概念来描述这一挑战，并呼吁在机构设计和激励机制上作出系统性应对\cite{petryk2023gharchive}。
本研究正是在这一宏观背景下，对开源激励机制中仍未被充分探讨的
信息维度作出回应。

\subsection{开发者参与动机的理论框架}

开源开发者动机研究的核心挑战，在于解释为何大量技能成熟的开发者
愿意将大量时间投入到没有直接经济回报的项目。
早期研究集中于内在动机的识别与刻画。
Shah（2006）对多个开源社区的质性研究发现，开发者参与的最初驱动力
往往是对技术问题本身的强烈兴趣——即所谓"抓自己的痒"，
开发者最初为解决自身实际技术需求而参与，贡献本身即是收益，
无需外部奖励来维持\cite{shah2006motivation}。
Davidson 等人（2014）进一步指出，利他主义——希望为社会留下有价值的成果——
在资深开发者中尤为突出，与兴趣、声誉追求共同构成多层次的内在动机结构\cite{davidson2014older}。

Lerner 和 Tirole（2002）则从经济学视角提出声誉机制假说：
在开源社区中，代码是公开的，每位开发者的贡献质量对所有人可见，
优秀贡献可以为开发者在同行社区中建立声誉，
而这一声誉可以转化为劳动力市场上的职业优势\cite{lerner2002simple}。
这一理论将声誉理解为一种"延迟货币"，
将表面上无偿的贡献解读为对人力资本的理性投资。
Fang 和 Neufeld（2009）的纵向研究进一步证实了职业发展动机的重要性，
发现将开源贡献视为职业发展手段的开发者，其持续参与的稳定性
显著高于纯粹基于兴趣的参与者\cite{fang2009understanding}。

在内在动机与外在动机的理论谱系中，
自我决定理论（Self-Determination Theory）提供了一个更为精细的分析框架。
该理论区分了不同程度的外在动机内化（internalization）程度：
当外在目标被个体充分内化——即个体认为追求该目标与其核心价值观一致时，
外在动机与内在动机不再对立，反而可以协同增强行为动机。
在开源社区中，当开发者将接受外部资助理解为"对自己的技术贡献价值的社会认可"
而非"为了金钱而工作"时，外在奖励并不必然触发挤出效应；
反之，若奖励被感知为一种外部控制手段，则内在动机的侵蚀效应更可能发生。
这一分析暗示，捐赠信息的披露方式——是以"认可贡献价值"还是以"购买服务"
的框架呈现——可能对其最终的激励效果产生关键影响\cite{benabou2003intrinsic}。

内在动机与外在动机并非简单相加，二者之间存在复杂的交互关系，
学界将其概括为"挤出效应（crowding-out effect）"。
Bénabou 和 Tirole（2003）在理论模型中论证，当个体内在动机足够强时，
引入外在奖励反而可能产生反效果：
外在奖励传递了"这项工作本身不够有价值"的隐性信号，
并且破坏了个体基于内在动机行动所获得的社会认可\cite{benabou2003intrinsic}。
Qiao 等人（2021）在在线志愿知识贡献平台的实证研究印证了这一机制，
发现货币激励在短期内确能提升贡献量，
但同时显著降低高内在动机用户的长期参与意愿\cite{qiao2021mitigating}。
此外，社会性动机同样不可忽视。
Yang 等人（2021）区分了社会影响对初始参与与持续参与的差异化效应：
社会认同和群体规范主要驱动开发者作出初始贡献，
而持续参与则更多依赖基于互动与信任建立起来的社区归属感\cite{yang2021differential}。
Sun 等人（2017）发现，货币奖励的效果受到社会连接性的显著调节——
在社交关系紧密的平台上，货币奖励与社会认可的结合
能够产生超越单一维度激励的协同效果\cite{sun2017motivation}。
在加密货币生态的特殊背景下，Kobayakawa 等人（2020）还发现
代币市场价格表现通过影响开发者对项目前景的预期，
间接影响了开发者的参与强度，说明经济信号在特定情境下
可以超越直接货币转移而发挥激励作用\cite{kobayakawa2020impact}。

\subsection{货币激励与捐赠赞助机制的研究综述}

尽管内在动机在开源参与中占据主导地位，随着开源软件战略重要性的持续上升，
外部货币激励机制也逐渐进入开源生态，并引发了学界的系统研究。
GitHub 于2019年推出 Sponsors 平台，允许个人和企业直接向开源开发者
支付定期赞助费用，为开源贡献提供了迄今最大规模的平台级经济支持基础设施。
Shimada 等人（2022）对该机制进行了最早的系统性研究，发现超过四分之三的
开发者在获得赞助后表现出贡献量的提升，且部分开发者将接受赞助本身
理解为对开源工作价值的认可，而非单纯的经济交换\cite{shimada2022github}。
Wang 等人（2022b）通过自然实验分析发现，赞助不仅正向影响受赞助开发者的贡献，
还会通过社区观察机制对周围开发者产生有限的溢出激励效应\cite{wang2022influence}。
Zhang 等人（2022）考察了"谁赞助谁"的决定因素，发现赞助机制
更倾向于奖励已在社区中建立声誉的活跃开发者，因而更多发挥
贡献后的事后认可功能，而非对潜在贡献的事前激励\cite{zhang2022who}。

然而，另一批研究呈现了截然不同的图景。
Overney 等人（2020）对多平台开源捐赠机制的大规模分析发现，
捐赠对开发者活动水平的整体提升作用十分有限，
多数项目即便获得捐赠，开发者的活动模式也并未发生显著变化\cite{overney2020not}。
Conti 等人（2023）追踪了 GitHub Sponsors 数年的纵向数据，
发现实际收到赞助的开发者在长期内表现出创新型贡献的减少，
经济激励可能将开发者的注意力从开放性技术探索
转移到能够维持赞助关系的稳定性工作上\cite{conti2023incentivizing}。
Wang 等人（2022a）在自然实验框架下发现，货币奖励产生了显著的
"知识溢出挤出"效应——货币化激励使贡献者减少了免费分享知识的行为\cite{wang2022monetary}。

GitHub Sponsors 研究与非盈利基金会资助研究之间还存在一个值得关注的结构性差异。
GitHub Sponsors 是一种"开发者对开发者"或"粉丝对开发者"的点对点直接支持机制，
赞助者可以自主选择受赞助对象，透明度和即时性均较高；
而以 Brink 为代表的专项基金会则扮演"中间人"角色，
通过集中募资后再经过申请审核程序将资金分配给具体开发者。
这种结构的差异带来了信息传播路径的根本不同：
在 GitHub Sponsors 模式下，赞助关系的建立本身是双方的直接互动，
赞助信号直达受赞助者，同时对旁观者具有一定的公开可见性；
而在基金会模式下，捐款信息的公开传播与资金的最终分配完全分离，
公告的受众远大于最终受益者，这一结构天然形成了可以独立考察信息效应的实验条件。
正是这种"信息先于货币传播、信息受众大于货币受益者"的特征，
使得基金会模式下的信息披露研究具有超越 GitHub Sponsors 研究的独特价值。

综合上述研究可以识别一个共同局限：
现有文献对货币激励的理解普遍聚焦于实际货币转移——某位开发者实际收到赞助金额这一事实，
而将激励效应等同于直接经济补偿的作用。
然而，捐赠公告所携带的信息价值——
在资金实际到达之前甚至之后，信息本身对未直接受益的开发者所可能产生的激励效应——
在所有这些研究中均未被单独识别和检验。
捐赠信息披露作为独立于货币转移的激励机制，其有效性究竟如何，
构成了本研究有别于既有文献的核心问题所在。

\section{信息披露的理论与实证}

\subsection{自愿信息披露理论}

信息披露研究的系统性发展始于会计学领域，并被管理学和信息系统学科广泛借鉴。
自愿信息披露理论（Voluntary Disclosure Theory）的核心命题是：
在信息不对称的环境中，拥有私有信息的组织会通过权衡披露成本与收益，
自主决定是否以及如何向外部受众传递信息。
当披露预期带来的收益——如降低信息使用者的不确定性、
提升对组织的信任、减少信息使用者要求的风险溢价——超过披露成本时，
组织将选择主动披露\cite{huang2021examining}。

对于非营利组织而言，信息披露具有维护公信力的特殊意义。
捐赠者和社会公众对非营利组织的资金来源与使用情况具有高度关注，
透明的财务信息披露是这类组织维持捐赠关系、建立社会信任的基础性手段。
对于 Bitcoin Brink 这类依赖社区捐赠运营的开源资助机构，
信息透明度不仅是合规性要求，更是其与开发者社区建立信任关系的核心机制。
Brink 在 X 平台公开每笔收到的捐款、在官网公布每位受资助开发者名单，
这种主动、详细的信息披露行为，既是对捐赠者问责的承诺，
也是向开发者社区传递"项目获得外部认可"这一积极信号的渠道。
从自愿披露理论的视角看，Brink 的披露行为本身即是一种精心设计的信息策略，
其对开发者行为的影响效果值得系统评估。

\subsection{信号理论}

信号理论（Signaling Theory）源自劳动经济学中 Spence（1973）
对劳动力市场信息问题的经典分析\cite{spence1973job}，
后被广泛应用于组织行为学、战略管理和信息系统领域。
信号理论的基本洞察是：在信息不对称的情境下，
处于信息劣势的一方（信号接收者）需要依靠可观察的信号
来推断其无法直接观察的信息；
而处于信息优势的一方（信号发送者）可以通过发送具有可信性的信号，
降低对方的不确定性并引导其作出有利的判断与行动。
一个有效信号需满足两个条件：
一是与所传递的内容真实相关（相关性），
二是具有足够的成本使得虚假信号不值得发送（可信性）\cite{connelly2011signaling}。

在开源捐赠信息披露的情境中，Brink 的两种披露构成了性质不同的信号。
非定向披露传递的是"项目获得社会认可与外部资金支持"的整体性信号，
类似于众筹平台上项目获得早期资助所产生的"热度信号"，
向整个开发者群体传递了"这个社区正在获得外界关注"的正面判断依据；
定向披露传递的则是"高质量贡献会获得认可与回报"的差异性信号，
通过明确揭示"谁因何贡献而获得资助"来强化努力—奖励的因果链条，
向未入选者传递了关于资助门槛和评判标准的具体信息，
从而影响其对"自己是否也有可能获得资助"这一问题的预期评估。
两类信号在内容、受众和激活机制上的本质差异，
构成了本研究分别检验其激励效应的逻辑起点。

\subsection{奖励信息披露与受众行为}

在信息披露对受众行为影响的实证研究中，众筹领域提供了与本研究最为接近的参照情境。
Huang 等人（2021）的元分析综合了多项众筹研究的结论，
发现项目进展公告（如资金筹集量的更新）能够显著提升后续的捐赠意愿，
即便公告内容并不伴随任何即时的经济回报——
观察者通过他人支持行为来修正对项目质量的判断，形成从众性的参与激励\cite{huang2021examining}。
这一机制说明，信息本身可以作为独立于货币的激励因素发挥作用，
其核心路径是通过改变受众对参与价值的预期评估来引导行为。

在知识贡献平台的研究中，Toubia 和 Stephen（2013）区分了
"内在效用"（来自贡献内容本身的满足）与"形象效用"
（来自贡献被他人看见和认可的满足），
发现形象效用在很大程度上依赖于平台所传递的"贡献受重视"的社会信号\cite{toubia2021attention}。
当平台通过公告、排行榜或特殊徽章等方式显式传递这类信号时，
用户的参与意愿会受到显著影响，尽管这些信号并未带来任何直接的金钱收益。
上述证据共同表明，信息披露能够通过激活受众期望、修正其对奖励可得性的判断，
产生独立且可测量的行为效应，为本研究的核心假设提供了实证支撑。

从比较视角看，Brink 的捐赠信息披露与上述研究情境存在若干重要的共性与差异。
共性在于，信息的公开传播先于或独立于实际经济转移而发生，
且信息受众（所有关注 Brink 的开发者）远大于实际经济受益者（获得资助的4名开发者）；
差异在于，开源开发者贡献的回报机制更为复杂——
他们面临的不仅是"是否参与"的二元决策，还有"投入多少"的连续决策，
且这种决策嵌入在多年的社区参与历史中，受到声誉、职业发展和内在动机的交织影响。
这意味着，Brink 信息披露的效应可能比众筹场景中的进展公告效应更为细微和情境化，
但其分析价值也在于：在这种高度复杂的动机结构中，
信息披露仍能产生可识别的增量效应，则意味着信息因素对开源开发者行为具有
超越噪音水平的独立影响力，这一发现将具有重要的理论意涵。

\section{期望理论与期望失验理论}

\subsection{期望理论}

期望理论（Expectancy Theory）由心理学家 Victor Vroom 于1964年在其工作动机研究中
正式提出\cite{vroom1964work}，是组织行为学和管理学领域被引用最广泛的动机理论之一。
该理论的核心命题是，个体的行为动机不由需求本身决定，
而由其对行动结果的认知预期决定：个体愿意付出努力的程度，
取决于三个认知因素的乘积。
其一是期望值（Expectancy），即个体相信自身的努力
能够达到预期绩效水平的主观概率；
其二是工具性（Instrumentality），即个体相信达到特定绩效
能够带来相应奖励的主观概率；
其三是效价（Valence），即个体对特定奖励结果的主观价值评估。
当任何一个因素趋近于零时，整体动机水平随之崩溃——
这解释了为什么在不信任奖励可得性或不重视奖励内容的环境中，
外部激励往往难以奏效。

期望理论最初用于解释雇佣关系中员工的努力水平，
但其"努力→绩效→奖励"的因果链逻辑具有高度的跨情境适用性，
被广泛延伸至志愿行为、公民参与和在线贡献等非雇佣情境。
在开源社区中，"奖励"的概念扩展至声誉提升、技术满足感和社区认可等多种形式；
但一旦外部货币激励机制进入社区（如 Brink 的资助项目），
货币奖励便成为开发者感知奖励结构的重要组成部分，
期望理论的逻辑随之获得了更直接的适用空间。

期望理论在在线贡献情境中的适用性已获得若干实证研究的初步支持。
在众筹领域，Huang 等人（2021）发现，项目进展信息更新
能够提升潜在支持者对项目最终成功的预期，进而提升其贡献意愿，
这一路径与期望理论的期望值维度高度吻合\cite{huang2021examining}。
在在线问答社区中，货币奖励机制（如平台积分或徽章）
被发现通过提升工具性感知（高质量回答确实会获得奖励）来激励知识贡献，
而非仅仅通过奖励的直接经济价值发挥作用。
这些研究共同指向一个结论：奖励信息的可见性与可信性——
即潜在贡献者是否能够清楚地看到并相信"努力→奖励"的实现路径——
对于外部激励能否有效作用于行为，具有与激励本身同等重要甚至更为关键的意义。
Brink 的信息披露实践，从信息可见性的角度为开发者提供了这种路径的具体依据，
其作用机制正是通过修正期望理论三要素的感知水平来实现的。

在本研究的具体情境中，Brink 的非定向披露通过向整个开发者社区传递
"资金存在、奖励可得"的信号，同时作用于期望理论的两个维度：
一方面，捐款金额的公开提升了开发者对资助项目规模与持续性的感知，
增强了奖励的效价（Valence）；
另一方面，资金的公开存在使开发者对"自己也有可能被资助"
产生了更积极的预期，提升了期望值（Expectancy）。
而定向披露通过揭示具体受资助者及其资助理由，
向未入选开发者传递了"优秀贡献确实会被认可和奖励"的证据，
强化了绩效与奖励之间的因果连接，提升了工具性（Instrumentality）。
基于这一分析，本研究预期两种披露均会正向影响开发者的贡献行为，
但其激活的认知路径有所不同。

\subsection{期望失验理论}

期望失验理论（Expectancy Disconfirmation Theory）由消费者行为研究者
Richard Oliver 于1980年提出，用于解释消费者满意度的形成机制\cite{oliver1980cognitive}。
该理论认为，个体满意度并非由实际结果的绝对水平决定，
而是由实际感知与事先期望之间的比较关系决定：
当实际结果超过期望时，个体产生正向失验（positive disconfirmation），
体验到满意乃至喜悦；
当实际结果低于期望时，个体产生负向失验（negative disconfirmation），
体验到失望乃至抱怨；
当结果与期望相符时，期望被确认，个体保持中性或满足状态。
这一框架最初在消费者满意度研究中取得广泛实证支持，
后被延伸应用于信息系统使用后评价\cite{abrate2021relationship}等多种情境，
并在解释重复使用意愿、持续贡献行为等方面展现出稳健的解释力。

期望失验理论的核心贡献在于揭示了期望水平作为比较基准的关键作用：
同样的实际结果，在不同期望背景下可能产生截然相反的情绪反应和后续行为。
对于未被 Brink 资助的开发者而言，
非定向披露所激活的"我也可能被资助"构成了一个比较基准；
当定向披露随后公布受资助者名单，这批开发者将自身状况与公告结果进行比较，
对于绝大多数未入选者而言，这一比较产生了明确的负向失验——
期望未能实现，由此产生的失望情绪可能削弱其参与意愿，
尤其是削弱其对后续非定向披露所激活期望的响应强度。
这种"泼冷水"效应，正是本研究预测定向披露对非定向披露具有负向调节效应的核心理论依据。

需要指出的是，期望失验效应的强度受到初始期望水平的调节。
根据 Abrate 等人（2021）对价格—体验不匹配的研究，
初始期望越高、结果的负向偏离越大，产生的失望情绪越强烈；
若初始期望本就较低，则负向失验的冲击也相应较弱\cite{abrate2021relationship}。
在本研究情境中，不同贡献水平的开发者对获得 Brink 资助的初始期望存在差异：
重度贡献者可能认为自己"更有资格"，因而初始期望更高，
负向失验后的动机削弱也更为明显；
轻度贡献者则可能将期望预设得较低，因而受定向披露影响相对有限。
这一理论预期为本研究后续的开发者异质性分析提供了重要指引。

\subsection{两种理论的整合：序贯披露的复合效应}

期望理论与期望失验理论在 Brink 两阶段披露的时间序列中形成了自然的互补关系，
共同描述了一条完整的认知动态链条：
非定向披露首先激活期望，在所有开发者中提升效价与期望值，初步促进贡献；
定向披露随后将这一过程推向结果阶段，其效应在两个层面同时展开——
对全体开发者而言，定向披露通过展示成功案例增强工具性，产生正向主效应；
而对大多数未入选者而言，定向披露又触发了负向期望失验，
削弱了其对后续非定向披露所激活期望的响应强度，产生负向调节效应。
这种"整体正向、交互负向"的复合结构，
在同一研究情境中为两种理论提供了同步检验的机会，
也是本研究在单一数据框架内整合两种理论的逻辑基础。

需要特别指出的是，期望理论与期望失验理论的整合在信息系统与在线平台研究中
已有先例，但具体应用场景各有侧重。
在消费者技术接受领域，Oliver（1980）的失验框架被 Bhattacherjee（2001）
等人扩展为IS持续使用研究的核心模型——"确认模型（Expectation-Confirmation Model）"，
认为用户在初次使用后会将实际体验与使用前期望进行比对，
确认或失验这一认知过程直接决定了其持续使用意愿。
尽管本研究的对象是开发者贡献行为而非技术产品的持续使用，
但其底层逻辑高度相似：非定向披露激活期望（类比技术使用前的期望形成），
定向披露公布受益人（类比使用后的实际结果感知），
期望与结果的落差（类比确认/失验）最终影响对后续披露的响应强度。
这种结构上的相似性不仅为本研究的理论借用提供了依据，
也意味着本研究的发现可与信息系统持续使用研究领域进行跨领域对话，
在理论贡献的辐射范围上具有一定的延展性。
此外，值得关注的是，期望水平本身在本研究情境中并不固定，
而是随着非定向披露次数的积累和社区对 Brink 运作模式的了解而动态调整。
随着开发者逐渐认识到获得 Brink 资助的实际门槛，
其初始期望的高估程度会随时间递减，负向失验效应也可能相应减弱。
这一动态性质提示，研究需要在面板模型中对时间趋势加以控制，
以避免将随时间演变的期望调整过程错误归因于特定披露事件的效应。

\section{情感与信息处理}

\subsection{情感即信息理论}

情感即信息理论（Affect-as-Information Theory）由 Schwarz 和 Clore（1983）
在一系列以情绪和判断为主题的经典实验研究中提出\cite{schwarz1983moods}。
该理论的核心主张是：个体在作出判断和决策时，
会将当前所感受到的情感状态作为一种关于处境质量的信息来源加以利用，
从而影响判断和行动的方向。
情感状态不仅仅是对事件的情绪反应，它本身携带着关于环境状态的认知内容：
积极情感传递"当前状况良好、无需采取额外行动"的信号，
促使个体倾向于维持现状或从事探索性活动；
负面情感则传递"当前存在问题或威胁、需要警惕与应对"的信号，
促使个体转向系统性分析和问题解决导向的行为。

在开源社区情境中，Issue 讨论区是社区成员集体情感最为集中体现的场所。
当项目遭遇技术瓶颈、安全漏洞或核心开发者之间的意见分歧时，
Issue 讨论的情感极性会出现明显的负面偏移——
语气趋于焦虑、沮丧或激烈，隐含着"项目面临困难、需要更多投入"的集体信号。
根据情感即信息理论，观察到这种集体负面情感的开发者
会将其解读为"项目处于需要行动的危机状态"的信号，
从而激发增加代码贡献（通过提交修复代码直接应对问题）
或增加讨论参与（通过协商寻求解决方案）的动机。
在本研究中，情感即信息理论构成了社区负面情感
正向影响代码活动和讨论活动这一预期（H4、H5）的核心理论依据。

值得注意的是，情感信号所触发的行为效应并非在情感发生的当期即刻呈现，
而是在情感信号被个体感知、加工后，经过一定的认知消化期才影响行为决策。
这意味着情感变化与行为变化之间存在一定的滞后关系，
这一预期与向量自回归模型所捕捉的动态传导机制天然契合，
使情感即信息理论框架与本研究所采用的时间序列方法论相互支撑。

\subsection{目标手段替代理论}

目标手段替代理论（Goal-Means Substitution Theory）源自目标系统理论
（Goal Systems Theory）的框架，由 Kruglanski 等人（2002）系统提出\cite{kruglanski2002theory}。
目标系统理论认为，目标与实现目标的手段之间形成关联网络：
当同一目标可通过多种手段实现时，若某一手段已被充分激活和使用，
与同一目标关联的其他替代手段的激活程度会相应降低。
其内在机制在于，目标激活（即个体对实现目标的需要感）驱动手段选择，
当某一手段的执行使目标的实现程度得到充分满足时，
对替代手段的需求感随之减弱；
加之个体时间与精力资源的有限性，已高度激活的手段会占用可用于替代手段的注意力资源，
进一步压缩其使用空间。

在开源社区协作的情境中，代码活动（提交代码修复问题）
与讨论活动（通过交流协商解决方案）可被理解为
以"解决项目问题"或"维持项目进展"为共同目标的两种替代性手段。
当社区成员已大量提交代码修复，有效推进了问题解决时，
"解决项目问题"这一目标通过代码手段得到了相当程度的满足，
对讨论这一替代手段的需求感相应降低。
因此，本研究预期代码活动的增加会在滞后一定时期后
对讨论活动产生负向的替代效应（H6）。
这一分析还暗示替代效应的强度具有情境依赖性：
在代码质量高、修复进展顺利的时期，替代效应更为明显；
而在代码进展受阻、问题悬而未决的时期，
讨论的必要性不会因代码活动增加而减少，替代效应随之减弱。

\subsection{开源社区中的情感研究综述}

随着自然语言处理技术的成熟，情感分析（Sentiment Analysis）
逐渐成为开源社区研究的重要工具。
Guzzi 等人（2013）对 Apache 和 Mozilla 项目邮件列表的分析发现，
开发者之间的情感沟通——表达感谢、沮丧、激动——
在技术决策讨论中占有相当比例，
这些情感内容构成了社区社会资本与协作质量的重要信号\cite{guzzi2013}。
Schroter 等人（2010）从问题报告者的情感表达切入，
发现情感较积极的 Bug 报告更可能获得开发者的响应与修复，
揭示了情感表达对技术协作结果的实际影响\cite{schroter2010}。
这些研究初步确立了情感在开源社区中扮演的信息传导角色，
但普遍采用截面分析或短期事件研究，难以捕捉情感变化与行为变化之间的动态时间关系。

更重要的局限在于，现有研究对情感与不同类型开发活动
（代码活动 vs. 讨论活动）之间差异化影响路径的关注不足，
更缺乏对代码活动与讨论活动之间动态替代关系的系统考察。
换言之，"情感如何分别影响代码贡献和讨论参与？
代码活动的增加是否会随后压缩讨论活动的空间？"
这一涉及情感—行为动态传导完整链条的问题，
目前在开源社区研究中尚属空白。

现有开源情感研究在方法论层面同样存在值得关注的局限。
早期研究大多使用基于词典的情感分析工具（如 SentiStrength、LIWC），
这类工具在通用文本上表现尚可，
但在高度技术化的开源讨论文本中，大量领域专属表达（如代码片段、版本号、
错误堆栈信息）会对情感评分产生系统性干扰，导致准确性受限。
近年来，基于预训练语言模型的情感分析方法（如 BERT 及其衍生模型）
在技术文本理解上取得了显著进步，但在开源社区情感分析中的系统性评估仍相对有限。
如何选择适合技术文本的情感分析工具、如何处理专业领域词汇对通用情感模型的干扰，
本身即是开源社区情感研究面临的方法论挑战之一，
本研究将在数据与方法章节对此作出专门说明，
并通过替换情感工具的稳健性检验验证主要结论对工具选择的不敏感性。
通过构建包含情感、代码活动和讨论活动的 VAR 模型，
本研究在方法论和研究问题两个层面均对既有研究作出实质性推进。

\section{开源社区中的活动类型与相互关系}

GitHub 作为全球最大的开源代码托管平台，
提供了高度结构化的开发活动记录，
这些活动记录构成了量化研究开源社区协作行为的核心数据资产\cite{petryk2023gharchive}。
在 GitHub 平台上，开发者可以参与多种形式的活动，
不同活动类型在与项目最终产出的直接关联程度、参与门槛以及
对外部信号的响应机制上均存在实质性差异，值得分别对待。

代码贡献活动（Code Activity）是开发者对项目代码库最直接的技术贡献，
包括 CommitCommentEvent（CCE）和 PullRequestEvent（PRE）两种形式。
代码贡献要求开发者具备足够的技术储备和对项目状态的深度了解，
通常需要投入相对连续且高度集中的时间，是反映开发者核心产出的最直接指标。
审查贡献活动（Review Activity），即 PullRequestReviewEvent（PRRE），
是开发者对他人提交代码进行质量评估的协作行为。
代码审查要求审查者深入理解项目规范和技术标准，
是维护代码质量、防止技术债务累积的关键机制，
在大型开源项目中也往往是新贡献能否被合并的决定性环节。
讨论活动（Discussion Activity），即 IssueEvent（IE），
涵盖问题报告、需求提案、功能讨论和技术争辩等多种形式，
参与门槛相对较低，任何具备基本技术背景的社区成员均可参与，
是社区协商和集体决策的主要渠道。

这三类活动不仅在功能上存在区别，
在对外部信息信号的响应机制上也可能呈现差异。
代码贡献所需时间成本最高，其变化往往反映深层的动机改变，
响应外部信号的时间窗口相对较长；
讨论活动时间成本较低，响应外部刺激的速度更快、弹性更大；
审查活动介于两者之间。
对三类活动分别建模，有助于揭示不同激励机制在不同活动层面的差异化影响，
从而对开源开发者行为结构作出更为细致的刻画\cite{schroter2010}。

此外，三类活动之间并非相互独立，而是通过开发者有限的时间与精力资源相互关联。
代码活动与讨论活动均服务于"推进项目解决问题"这一共同目标，
但所需资源和关注焦点不同：代码活动要求高度专注的技术投入，
而讨论活动更依赖对社区动态的关注和沟通意愿。
本研究引入目标手段替代理论，对代码活动与讨论活动之间可能存在的替代关系
提出明确的理论预期，并通过 VAR 模型对这一动态路径进行实证检验，
从而将活动类型的理论区分与动态因果分析有机结合起来。

从研究设计的角度看，三类活动的分开测量还具有规避"合并谬误"的方法论意义。
如果将所有开发活动合并为单一指标，
可能会掩盖激励信号对不同活动类型的差异化效应——
例如，非定向披露可能更强烈地激励需要较低门槛的讨论参与，
而对需要大量时间投入的代码贡献影响有限；
定向披露所展示的榜样效应则可能更强烈地激励能够直接展示技术能力的代码贡献。
类似地，社区情感对讨论活动的即时影响可能强于对代码活动的影响，
因为讨论对情感信号的响应更为迅速，而代码提交往往需要较长的准备周期。
通过对三类活动分别建模，本研究能够更精细地捕捉激励机制的作用边界，
避免因指标聚合而遮蔽有价值的异质性信息。
这种细粒度的活动分类，也使本研究的发现对开源社区的实践管理
具有更具体的参考价值——管理者可以据此判断哪种类型的激励策略
对激活哪种类型的贡献活动最为有效。

\section{本章小结}

本章围绕两项研究的理论需求，系统梳理了五个相互关联的理论与文献板块。
开源开发者激励文献揭示了内在动机主导下货币激励可能产生挤出效应的复杂格局，
并指出现有研究忽视了信息披露本身的激励价值；
自愿信息披露理论和信号理论为理解 Brink 两阶段披露机制的信号性质提供了概念工具；
期望理论和期望失验理论为预测两种披露的主效应及其交互效应
提供了心理机制层面的解释框架；
情感即信息理论和目标手段替代理论为社区情感影响开发活动、
并进而触发代码-讨论替代关系提供了理论依据；
对开源社区活动类型及其相互关系的梳理，
确立了本研究区分三类贡献活动分别建模的方法论合理性。
以上理论框架共同构成了后续两项实证研究的逻辑基础，
并在各章假设推导与结果解读中得到具体应用。
图~\ref{fig:theory-framework} 以可视化方式呈现了上述各理论模块
与研究一、研究二之间的对应关系，
有助于读者在进入实证章节之前把握本文整体的理论结构。

\begin{figure}[htbp]
    \centering
    \fbox{\parbox{0.9\textwidth}{\centering
        \vspace{0.8cm}
        \textbf{研究一（捐赠信息披露）}\\[0.4em]
        自愿信息披露理论 $+$ 信号理论
        $\longrightarrow$ 两种披露的信号性质 \\[0.3em]
        期望理论
        $\longrightarrow$ H1（非定向披露正向主效应）$+$ H2（定向披露正向主效应）\\[0.3em]
        期望失验理论
        $\longrightarrow$ H3（定向披露对非定向披露的负向调节）\\[0.8em]
        \textbf{研究二（社区情感与开发活动）}\\[0.4em]
        情感即信息理论
        $\longrightarrow$ H4（负面情感 $\to$ 代码活动）$+$ H5（负面情感 $\to$ 讨论活动）\\[0.3em]
        目标手段替代理论
        $\longrightarrow$ H6（代码活动 $\to$ 讨论活动的替代效应）
        \vspace{0.8cm}
    }}
    \caption{理论框架与研究假设的对应关系}
    \label{fig:theory-framework}
\end{figure}


%%====================================================================
\chapter{研究设计与数据概览}
%%====================================================================

本章作为两项实证研究的共同基础，统一呈现整体研究框架、数据来源与融合流程、
方法论选择逻辑及伦理考量。通过在独立章节中整合上述共性内容，
可以避免两项研究各自孤立叙述所造成的论文整体感不足，
同时为读者在进入具体实证章节之前建立清晰的全局认知。

\section{整体研究框架}

本文围绕"比特币开源社区开发者活动的影响因素"这一核心命题，
设计了两项在研究对象、数据基础和时间窗口上相互衔接、
在研究问题和方法论上相互补充的实证研究。
图~\ref{fig:overall-framework} 呈现了两项研究在整体逻辑中的位置与关系。

\begin{figure}[htbp]
    \centering
    \fbox{\parbox{0.9\textwidth}{
        \centering
        \vspace{0.6cm}
        \textbf{研究情境：比特币 GitHub 开源社区（2022年1月—2023年12月）}\\[0.8em]
        \begin{tabular}{p{0.38\textwidth}p{0.04\textwidth}p{0.38\textwidth}}
        \centering\textbf{研究一} & & \centering\textbf{研究二} \\[0.2em]
        \centering 外部信息输入 & & \centering 内部信息积累 \\
        \centering Brink 捐赠信息披露 & & \centering 社区 Issue 讨论情感 \\[0.4em]
        \centering $\downarrow$ & & \centering $\downarrow$ \\[0.2em]
        \centering 个体开发者活动（周级面板） & & \centering 社区聚合活动（周级时序） \\[0.4em]
        \centering 面板固定效应 + 2SLS & & \centering VAR + 格兰杰因果检验 \\
        \end{tabular}
        \vspace{0.6cm}
    }}
    \caption{整体研究框架}
    \label{fig:overall-framework}
\end{figure}

研究一以\textbf{个体开发者}为分析单元，考察外部非盈利基金会（Bitcoin Brink）
两种捐赠信息披露事件对开发者活动水平的因果影响。
研究设计的核心优势在于：通过剔除实际受资助的4名开发者，
将分析对象限定在未直接受益于捐款的1,613名开发者，
从而实现信息效应与货币效应的干净分离。
固定效应模型通过控制开发者个体特征和时间共同趋势，识别信息披露的净效应；
工具变量方法进一步解决捐款金额与社区活动之间可能存在的反向因果问题。

研究二以\textbf{社区整体}为分析单元，考察比特币 Issue 讨论区中提取的情感极性指标
如何在时间维度上影响代码活动与讨论活动，
以及两类开发活动之间是否存在动态的替代关系。
研究二将情感、代码活动和讨论活动三个变量均视为内生变量，
通过向量自回归（VAR）模型捕捉多变量时间序列系统中的动态反馈机制，
并借助格兰杰因果检验和脉冲响应函数对影响路径进行识别与可视化。

两项研究在研究对象上有所不同——研究一关注个体层面的动机响应，
研究二关注社区层面的集体动态——但共享相同的时间窗口（2022—2023年）、
相同的核心活动变量（CCE、PRE、PRRE、IE）以及相同的底层数据来源（GH Archive）。
这种"同一情境、互补视角、多层分析"的设计，使两项研究能够从不同维度共同支撑
"信息层面因素深刻影响开发者行为"这一核心论点。

\section{数据来源总览与融合流程}

\subsection{全部数据源汇总}

本研究共整合六类数据源，覆盖开发者行为、信息披露事件、
市场环境和社会关注度等维度。
表~\ref{tab:data-overview} 对全部数据源的基本信息进行汇总说明。

\begin{table}[htbp]
    \centering
    \caption{全部数据来源汇总}
    \label{tab:data-overview}
    \begin{tabularx}{\textwidth}{lXll}
        \toprule
        数据类型 & 数据源与采集方法 & 时间粒度 & 用途 \\
        \midrule
        GitHub 开发活动
            & GH Archive（gharchive.org）；BigQuery 批量下载
            & 事件级→周级聚合
            & 研究一、二因变量 \\
        非定向披露信息
            & X 平台 @brink\_bitcoin 历史推文；X API v2
            & 事件级
            & 研究一核心自变量 \\
        定向披露信息
            & brink.dev 官网存档；手工整理 + Wayback Machine
            & 事件级
            & 研究一核心自变量 \\
        Bitcoin 市场数据
            & Yahoo Finance；yfinance API
            & 日级→周级
            & 研究一控制变量 \\
        Google Trends 搜索指数
            & Google Trends；pytrends 接口，美国地区
            & 周级
            & 研究一工具变量 \\
        Issue 评论文本
            & GH Archive + GitHub REST API；关键词过滤 + 分页拉取
            & 评论级→周级聚合
            & 研究二情感分析 \\
        \bottomrule
    \end{tabularx}
\end{table}

GH Archive 是由 Ilya Grigorik 维护的 GitHub 事件归档服务，
自2011年起持续以小时为粒度记录 GitHub 平台上的全量公开事件，
并通过 Google BigQuery 提供免费的大规模查询接口。
本研究从 GH Archive 中提取 bitcoin/bitcoin 主仓库及其关联仓库的四类事件：
CommitCommentEvent（CCE）、PullRequestEvent（PRE）、
PullRequestReviewEvent（PRRE）和 IssueEvent（IE），
分别对应代码注释贡献、代码提交与合并请求、代码审查和议题讨论四种核心活动类型。

Google Trends 数据通过 pytrends Python 接口采集，
以关键词"Brink Bitcoin"在美国地区的周搜索量指数（0—100标准化分值）
衡量公众对 Brink 基金会的关注热度。
选择美国地区的理由在于，比特币核心开发者及主要捐赠方以美国为主要分布地区，
美国地区的搜索热度对捐款规模具有更强的预测能力，
在第一阶段回归中与内生变量的相关性更为突出。

\subsection{数据清洗与匹配流程}

鉴于六类数据源在格式、粒度和标识符体系上存在差异，
需要经历系统的清洗与匹配流程方可形成可分析的面板数据集。
本研究按照以下五个步骤依序完成数据整合：

\textbf{第一步，开发者 ID 统一化与机器人过滤。}
GH Archive 事件中的行为主体以 actor\_id（数字型）和 actor\_login（字符串型）
双字段记录。本研究将两者双向映射，以确保同一开发者在不同时期使用不同账号名时
不被误计为多人。机器人账号通过账号名称正则过滤剔除，
筛除规则包括名称包含"bot"、"[bot]"、"robot"等标记词的账号，
以及 GitHub Actions、Dependabot 等平台内建自动化账号，
此步骤共识别并剔除约30个机器人账号。

\textbf{第二步，时间对齐。}
所有数据源统一以协调世界时（UTC）为时区基准，
以 ISO 8601 周标准（周一为一周起始日）为时间颗粒，
将全部事件时间戳转换为对应 ISO 周标识符（格式：YYYY-WXX），
确保六类数据可以在相同的时间键上进行关联合并。

\textbf{第三步，样本筛选。}
保留2022年1月（2022-W01）至2023年12月（2023-W52）共105个 ISO 周内，
在 bitcoin/bitcoin 主仓库有至少一次活动记录的开发者。
进一步剔除在观测期内实际获得 Bitcoin Brink 全职资助的4名开发者，
确保研究一的最终样本中所有个体均未因捐赠而直接受益，
保证回归系数反映的是信息披露的纯信息效应。
最终样本包含1,613名开发者，共约169,365个开发者-周观测单元。

\textbf{第四步，披露事件标记。}
将 Brink 在 X 平台发布的非定向披露公告日期和在官方网站发布的定向披露公告日期，
分别精确映射至对应 ISO 周，生成非定向披露金额变量
（披露金额取自然对数，无披露周取0）和定向披露虚拟变量
（定向披露公布后各周取1，公布前取0）。

\textbf{第五步，情感聚合。}
对 GH Archive 中 bitcoin/bitcoin 主仓库 Issue 区的评论文本，
逐条进行情感分析，获得每条评论的情感极性得分（compound score，
取值区间$[-1, 1]$，正值代表正面情感，负值代表负面情感）。
在此基础上，取每周内所有评论情感得分的算术均值，
作为当周社区整体情感极性指标，供研究二的 VAR 分析使用。

\section{方法论总述}

\subsection{两种方法的问题定位}

本文两项研究面临性质不同的计量问题，因而选择了互补的方法论路径。
表~\ref{tab:method-comparison} 对两种方法的核心特征与适用场景进行对比说明。

\begin{table}[htbp]
    \centering
    \caption{两种方法的问题定位与适用场景对比}
    \label{tab:method-comparison}
    \begin{tabularx}{\textwidth}{lXXX}
        \toprule
        方法 & 适用问题类型 & 应用于 & 核心解决的挑战 \\
        \midrule
        面板固定效应 + 2SLS
            & 识别外生冲击对个体的因果效应
            & 研究一
            & 个体异质性 + 内生性（反向因果）\\
        向量自回归（VAR）
            & 刻画多变量系统的动态互动机制
            & 研究二
            & 双向因果（三变量均内生，无先验主从关系）\\
        \bottomrule
    \end{tabularx}
\end{table}

研究一的核心任务是识别捐赠信息披露对开发者活动的因果效应。
这一场景面临两类典型内生性威胁：
其一是个体异质性——不同开发者在技能水平、内在动机和活跃程度上
存在持续性的横截面差异，若不加控制，将产生遗漏变量偏误；
其二是反向因果——社区整体的活跃程度可能影响外部捐赠者的捐赠意愿，
进而影响 Brink 的披露频率与金额，导致因果方向混淆。
面板固定效应模型可以有效处理第一类威胁，
而工具变量方法则进一步解决第二类威胁，
两者的结合构成了识别纯信息效应的核心识别策略。

研究二的核心任务是描述情感、代码活动和讨论活动三个变量之间的动态反馈机制。
与研究一不同，研究二并不存在明确的"外生冲击者"与"被冲击者"的先验区分——
情感可能影响活动，活动也可能影响情感；代码活动可能影响讨论活动，反之亦然。
向量自回归（VAR）模型将三个变量均视为内生变量，
通过联立方程组同时捕捉所有方向的动态反馈关系，
无需对因果方向作出不可检验的先验假定。

\subsection{固定效应模型的识别逻辑}

固定效应模型通过在基准回归方程中引入两维固定效应来控制两类系统性偏误。
开发者个体固定效应 $u_i$ 吸收了所有不随时间变化的个体特征——
包括编程技能水平、在比特币社区的参与年限、初始内在动机偏好，
以及其他可能同时影响活动水平和对披露信号响应强度的稳定个体属性。
时间固定效应 $\theta_t$ 则吸收了所有开发者在同一周共同面临的宏观趋势——
比特币价格的长期周期、加密货币监管政策的阶段性变化、
整个比特币生态的技术演进节奏，以及其他在各周内对所有开发者产生同质影响的外部因素。

在双维固定效应的控制下，核心系数 $\beta_1$ 所捕捉的，
是"在排除开发者个体差异和全局时间趋势之后，
捐赠披露信息的变动对个体开发者活动水平的净效应"，即信息披露的纯增量影响。
在此基础上，工具变量方法通过引入"Brink Bitcoin"Google Trends 搜索量指数
作为捐款金额的工具变量，进一步解决捐款金额变量本身的内生性问题。
有效工具变量需满足两个条件：
相关性条件要求工具变量与内生变量显著相关——本研究中公众的 Google 搜索热度
与 Brink 收到的捐款金额存在正向相关关系，
第一阶段回归的 F 统计量超过10的传统强工具变量门槛，确认相关性成立；
排他性条件要求工具变量仅通过内生变量影响因变量，不存在直接的影响路径——
公众在 Google 上搜索"Brink Bitcoin"的行为本质上反映的是社会对该机构的信息兴趣，
不直接决定比特币 GitHub 上某位特定开发者当周的贡献行为，满足排他性约束的理论论据。

\subsection{VAR 模型的识别逻辑}

向量自回归模型将情感（$\text{Sentiment}_t$）、代码活动（$\text{CodeActivity}_t$）
和讨论活动（$\text{DiscussActivity}_t$）三个变量组成系统向量 $\mathbf{y}_t$，
通过联立方程形式捕捉变量间的动态反馈：

\begin{equation}
\label{eq:var-overview}
\mathbf{y}_t = \mathbf{c} + \sum_{k=1}^{p} \mathbf{A}_k \mathbf{y}_{t-k} + \boldsymbol{\varepsilon}_t
\end{equation}

其中 $\mathbf{A}_k$ 为第 $k$ 期的系数矩阵，$\boldsymbol{\varepsilon}_t$ 为白噪声残差向量。
滞后阶数 $p$ 依据 AIC/BIC 信息准则确定，
优先选择使信息准则最小化的较小滞后阶数，
以在模型拟合与过拟合风险之间取得平衡。

格兰杰因果检验基于受限与无约束 VAR 模型的似然比统计量，
在统计层面识别"变量 X 的历史值是否有助于预测变量 Y 的当前值"——
若有助于预测，则称 X 格兰杰引起（Granger-cause）Y。
需要指出的是，格兰杰因果的成立是统计层面的预测能力判据，
并不等同于结构性因果关系，但结合本研究的理论预期和情境合理性，
可以作为动态影响路径存在的重要证据。

脉冲响应函数（IRF）采用 Cholesky 分解对冲击进行正交化处理，
排列顺序为"情感 → 代码活动 → 讨论活动"。
这一顺序的理论依据在于：情感状态作为信息环境的即时反映，
在逻辑上先于行为响应形成；代码活动由于时间成本较高，
对情感信号的响应速度相对迟缓；
而讨论活动参与门槛较低，对外部信号的即时响应更为灵敏，
因此在三者的因果顺序中居于末位。方差分解（FEVD）进一步量化每个变量
对系统中其他变量预测误差方差的贡献比例，
有助于评估各变量在整体动态系统中的相对重要性。

\section{伦理考量}

\subsection{数据公开性与合规性}

本研究所使用的全部数据均来源于公开渠道，研究性使用符合相关平台服务条款及
学术研究惯例。GH Archive 是 GitHub 对外开放的事件归档服务，
采用 CC BY 4.0 授权协议，允许研究性使用与再发布。
X 平台 @brink\_bitcoin 账号的推文内容为公开账号的主动发布信息，
不涉及私信或受访问控制的内容。
Brink 官网公布的受资助开发者名单是非盈利组织主动向社会公开的资助记录，
属于自愿信息披露行为，本研究对其进行学术性整理和引用具有充分的合理性。
Google Trends 和 Yahoo Finance 提供的数据均为面向公众的聚合统计数据，
不包含个人识别信息，研究使用不存在合规障碍。

\subsection{开发者隐私保护}

本研究在数据处理和结果呈现中充分考虑了开发者的隐私保护。
分析所使用的开发者标识符为 GitHub 平台的公开用户名（actor\_login）
及数字型 ID（actor\_id），不涉及开发者的私人联系方式、地理位置信息或
任何未经公开披露的个人数据。
在汇报统计结果时，所有分析均以聚合形式呈现——
回归系数反映的是样本整体的平均处理效应，
不对特定开发者的个体行为作出点名式描述。
在附录的披露事件列表中，出现的受资助开发者姓名
均来自 Brink 官方网站已公开发布的信息，未超出已公开范围。

\subsection{数据安全与可复现性}

原始数据存储于本地加密硬盘及受访问控制的服务器环境中，
研究过程中不对原始数据进行外部分发或共享。
数据分析所使用的 Python 脚本和 R 代码将在论文答辩完成后，
通过匿名化的 GitHub 代码仓库对外开放，以支持研究结果的可复现性审查。
代码仓库将不包含任何可识别作者信息的内容，
以满足同行评审过程中的匿名性要求。


%%====================================================================
\chapter{研究一：捐赠信息披露对开发者活动的影响}
%%====================================================================

\section{研究场景：Bitcoin Brink 生态}

\subsection{Brink 基金会简介}

Bitcoin Brink（以下简称"Brink"）成立于2020年，
是目前规模最大的专注于比特币核心开发的非盈利资助机构。
Brink 采用众筹模式，面向个人和机构收集捐款，
并通过严格的申请审核流程遴选受资助开发者，
为其提供全职从事比特币开发的资金支持。

Brink 的运营在信息披露层面形成了清晰可观测的两个节点（图~\ref{fig:brink-flow}），
本研究正是以此为核心研究场景。

\begin{figure}[htbp]
    \centering
    % \includegraphics[width=0.85\textwidth]{figures/brink-flow.pdf}
    \fbox{\parbox{0.85\textwidth}{\centering
        \vspace{1cm}
        [图：Brink 运营流程示意图]\\
        社会捐赠 → Brink 收款 → \textbf{X 平台公告（非定向披露）} \\
        \quad\quad\quad\quad\quad\quad\quad\quad\quad\quad ↓ \\
        开发者申请 → 审核 → \textbf{官网公布受资助名单（定向披露）}
        \vspace{1cm}
    }}
    \caption{Brink 捐赠信息披露流程}
    \label{fig:brink-flow}
\end{figure}

\subsection{非定向披露与定向披露的区别}

\textbf{非定向披露（Untargeted Disclosure）：}
当 Brink 收到捐款后，会在其 X 平台官方账号发布公告，
感谢捐赠方并公布捐款金额，但不说明资金将资助哪位具体开发者。
此类公告面向所有关注者，信息具有普惠性——
每位开发者均可将自己视为潜在的受益者。

\textbf{定向披露（Targeted Disclosure）：}
经过严格审核流程后，Brink 在其官方网站发布受资助开发者名单，
明确受资助者的姓名及资助期限。
此类披露揭示了"谁"获得了资金，具有明显的指向性和排他性。

在本研究观测期（2022年1月—2023年12月）内，
Brink 共发布非定向披露事件19次（其中15次含具体金额）；
最大规模的定向披露发生于2022年4月4日（共资助4名开发者）。

\subsection{样本设计的优越性}

本研究从观测样本中剔除了4名实际受到资助的开发者，
以确保研究对象均为未获实际资金的开发者，
从而将分析聚焦于信息披露的纯信息效应，
排除了实际货币激励对结果的干扰。
最终样本包含1613名开发者，105周，共约169,365个开发者-周观测。

\section{研究假设}

\subsection{非定向披露的正向主效应（H1）}

根据期望理论，非定向捐赠披露向所有开发者传达了奖励存在、
奖励可得的积极信号，从而提升了效价与期望值，激励贡献行为。

\begin{description}
    \item[H1] 非定向捐赠披露正向影响开发者的整体贡献活动。
    \item[H1a] 非定向捐赠披露正向影响开发者的代码贡献。
    \item[H1b] 非定向捐赠披露正向影响开发者的审查贡献。
    \item[H1c] 非定向捐赠披露正向影响开发者的讨论活动。
\end{description}

\subsection{定向披露的正向主效应（H2）}

根据期望理论的工具性维度，定向披露通过展示"绩效→奖励"的因果链，
使非受助开发者相信自身努力同样可以获得认可，激励其增加贡献。

\begin{description}
    \item[H2] 定向捐赠披露正向影响开发者的整体贡献活动。
    \item[H2a] 定向捐赠披露正向影响开发者的代码贡献。
    \item[H2b] 定向捐赠披露正向影响开发者的审查贡献。
    \item[H2c] 定向捐赠披露正向影响开发者的讨论活动。
\end{description}

\subsection{定向披露对非定向披露效应的负向调节（H3）}

根据期望失验理论，非定向披露在开发者中激活了"可能成为受益者"的期望；
定向披露公布后，绝大多数未入选者经历负向期望失验，
失望情绪削弱了其对后续非定向披露的响应强度。
因此，定向披露的出现将负向调节非定向披露的激励效果。

\begin{description}
    \item[H3] 定向披露负向调节非定向披露对开发者贡献的正向效应。
    \item[H3a] 该负向调节效应在代码贡献上成立。
    \item[H3b] 该负向调节效应在审查贡献上成立。
    \item[H3c] 该负向调节效应在讨论活动上成立。
\end{description}

研究模型如图~\ref{fig:research-model-1} 所示。

\begin{figure}[htbp]
    \centering
    \fbox{\parbox{0.85\textwidth}{\centering
        \vspace{1cm}
        [图：研究一假设模型]\\
        非定向披露 $\xrightarrow{+}$ 开发者贡献 \\
        定向披露 $\xrightarrow{+}$ 开发者贡献 \\
        定向披露 $\times$ 非定向披露 $\xrightarrow{-}$ 开发者贡献
        \vspace{1cm}
    }}
    \caption{研究一假设模型}
    \label{fig:research-model-1}
\end{figure}

\section{数据与变量}

\subsection{数据来源}

\begin{table}[htbp]
    \centering
    \caption{研究一数据来源}
    \label{tab:data-source-1}
    \begin{tabularx}{\textwidth}{lXl}
        \toprule
        数据类型 & 来源与说明 & 粒度 \\
        \midrule
        GitHub 开发活动
            & GH Archive（\url{https://www.gharchive.org}），
              涵盖 CCE、PRE、PRRE、IE 四类事件
            & 开发者—周 \\
        非定向披露事件
            & X 平台 @brink\_bitcoin 账号，共19次披露，
              15次含具体金额
            & 周 \\
        定向披露事件
            & Brink 官网（\url{https://brink.dev}），
              受资助名单及公告日期
            & 事件 \\
        Bitcoin 市场数据
            & Yahoo Finance，周收益率和交易量
            & 周 \\
        社会关注度
            & Google Trends，关键词"Brink Bitcoin"搜索指数
            & 周 \\
        \bottomrule
    \end{tabularx}
\end{table}

\subsection{关键变量定义}

\begin{table}[htbp]
    \centering
    \caption{研究一关键变量定义}
    \label{tab:variables-1}
    \begin{tabularx}{\textwidth}{llX}
        \toprule
        变量 & 类型 & 定义 \\
        \midrule
        $\text{develop}_{it}$
            & 因变量
            & 开发者 $i$ 在第 $t$ 周的代码贡献活动数量，
              $\text{develop}_{it} = \text{CCE}_{it} + \text{PRE}_{it}$ \\
        $\text{review}_{it}$
            & 因变量
            & 开发者 $i$ 在第 $t$ 周的审查贡献数量，
              $\text{review}_{it} = \text{PRRE}_{it}$ \\
        $\text{issue}_{it}$
            & 因变量
            & 开发者 $i$ 在第 $t$ 周的讨论活动数量，
              $\text{issue}_{it} = \text{IE}_{it}$ \\
        $\ln(\text{DonationAmount}_t)$
            & 核心自变量
            & 第 $t$ 周 Brink 在 X 平台公布的捐款金额（美元）的对数，
              度量非定向披露强度 \\
        $\text{TargetedDisclosure}_t$
            & 核心自变量
            & 虚拟变量，第 $t$ 周起 Brink 已披露受资助开发者名单取1，否则取0 \\
        $\text{GoogleSearchIndex}_t$
            & 工具变量
            & 第 $t$ 周"Brink Bitcoin"在 Google Trends 上的搜索量指数 \\
        \bottomrule
    \end{tabularx}
\end{table}

\subsection{描述性统计}

\todo{插入描述性统计表（Table 4 from sample paper）}

\section{计量模型}

\subsection{基准固定效应模型}

本研究采用个体—周两维固定效应模型，
同时控制个体异质性（$u_i$）和共同时间趋势（$\theta_t$）：

\begin{equation}
\label{eq:fe-model}
\begin{aligned}
\ln(\text{Contribution}_{it}) &= \beta_1 \ln(\text{DonationAmount}_t) \\
    &\quad + \beta_2 \text{TargetedDisclosure}_t \\
    &\quad + \beta_3 \ln(\text{DonationAmount}_t) \times \text{TargetedDisclosure}_t \\
    &\quad + \boldsymbol{\gamma}' \mathbf{X}_{it} + u_i + \theta_t + \varepsilon_{it}
\end{aligned}
\end{equation}

其中，$\mathbf{X}_{it}$ 为控制变量向量，包括：
滞后一期因变量、Bitcoin 周收益率、周交易量、
开发者发起的 Fork 数量和 Watch 数量。

\subsection{工具变量方法（2SLS）}

\textbf{内生性来源：}方程 \eqref{eq:fe-model} 中存在潜在的反向因果问题——
社区开发活动水平可能影响外部捐赠规模，进而影响 Brink 的披露行为。
为解决此内生性问题，本研究采用工具变量方法。

\textbf{工具变量：}以"Brink Bitcoin"的 Google Trends 搜索量指数
作为 $\ln(\text{DonationAmount}_t)$ 的工具变量。
\begin{itemize}
    \item \textbf{相关性}：公众对 Brink 的搜索关注度与其收到的捐款金额正相关——
          关注度高时往往也是捐款较多的时期；
    \item \textbf{排他性}：搜索指数反映的是社会公众的信息搜寻行为，
          不直接决定特定比特币开发者的个体贡献水平，
          满足工具变量排他性约束。
\end{itemize}

两阶段最小二乘估计（2SLS）模型：
\begin{align}
\text{（第一阶段）}\quad
    \ln(\widehat{\text{DonationAmount}}_t)
        &= \alpha_0 + \alpha_1 \text{GoogleSearchIndex}_t
           + \boldsymbol{\alpha}' \mathbf{X}_{it} + \eta_{it}
\label{eq:first-stage} \\
\text{（第二阶段）}\quad
    \ln(\text{Contribution}_{it})
        &= \beta_1 \ln(\widehat{\text{DonationAmount}}_t)
           + \beta_2 \text{TargetedDisclosure}_t \notag \\
        &\quad + \beta_3 \ln(\widehat{\text{DonationAmount}}_t)
           \times \text{TargetedDisclosure}_t \notag \\
        &\quad + \boldsymbol{\gamma}' \mathbf{X}_{it}
           + u_i + \theta_t + \varepsilon_{it}
\label{eq:second-stage}
\end{align}

\section{实证结果}

\subsection{主效应结果}

\todo{插入主回归结果表（对应 sample paper 中 Table 5）}

主要发现如下：
\begin{itemize}
    \item $\ln(\text{DonationAmount}_t)$ 的系数显著为正（代码贡献：0.001, $p < 0.05$；
          审查贡献：0.002, $p < 0.01$），\textbf{H1 成立}；
    \item $\text{TargetedDisclosure}_t$ 的系数显著为正（审查贡献：0.013, $p < 0.01$），
          \textbf{H2 成立}；
    \item 交互项系数显著为负（代码贡献：$-0.000$, $p < 0.05$；
          审查贡献：$-0.001$, $p < 0.01$），\textbf{H3 成立}。
\end{itemize}

\subsection{稳健性检验}

\subsubsection{不同时间窗口的事件研究}

\todo{事件研究（Event Study）结果}

\subsubsection{替换自变量度量方式}

\todo{使用虚拟变量 donationOrNot 替换连续金额变量后的结果}

\section{机制分析}

\subsection{开发者贡献量的异质性分析}

基于开发者在观测期内的历史总贡献量，将样本分为
\textbf{重度贡献者}（Heavy Contributors）与\textbf{轻度贡献者}（Light Contributors）两组。
\todo{分组回归结果及解释}

\subsection{开发者项目参与范围的异质性分析}

基于开发者参与的 GitHub 仓库数量，将样本分为
\textbf{专注型}（参与仓库数较少）与\textbf{分散型}（参与仓库数较多）两组。
\todo{分组回归结果及解释}

\subsection{开发者知名度的异质性分析}

以 GitHub Follower 数量的中位数为界，
将样本分为\textbf{高知名度开发者}（前50\%）和\textbf{低知名度开发者}（后50\%）。
\todo{分组 2SLS 结果及解释}

\section{小结}

本章基于期望理论与期望失验理论，
实证检验了 Bitcoin Brink 基金会两种捐赠信息披露方式对开发者活动的影响。
研究结果表明：两种披露均能正向激励开发者贡献，
但定向披露通过引发期望失验而削弱了非定向披露的长期激励效果。
这一发现强调了捐赠信息披露设计中"分寸感"的重要性，
并揭示了信息本身作为独立激励因素的有效性。


%%====================================================================
\chapter{研究二：社区讨论情感对开发者活动的影响}
%%====================================================================

\section{研究背景与问题}

开源社区的 Issue 讨论区是开发者交流技术问题、协商解决方案的主要场所，
蕴含丰富的社区情感信息——从对技术突破的兴奋，到对 Bug 的挫败，
再到代码审查中的意见分歧，情感极性随时间动态波动。

本章聚焦以下问题：
社区讨论情感如何在时间维度上影响开发者的代码活动与讨论活动？
代码活动的增加是否会对讨论活动产生替代效应？
这些影响路径如何从理论层面加以解释？

\section{理论假设}

基于情感即信息理论与目标手段替代理论，本章提出以下假设：

\begin{description}
    \item[H4] 滞后一期的负面情感（情感极性下降）正向影响代码活动。
    \item[H5] 滞后一期的负面情感（情感极性下降）正向影响讨论活动。
    \item[H6] 滞后两期的代码活动增加负向影响讨论活动。
\end{description}

完整影响路径如图~\ref{fig:path-model-2} 所示：

\begin{figure}[htbp]
    \centering
    \fbox{\parbox{0.85\textwidth}{\centering
        \vspace{0.5cm}
        \textbf{情感（$t-1$）} $\xrightarrow{-}$
        \textbf{代码活动（$t$）} $\xrightarrow{-}$
        \textbf{讨论活动（$t+1$）} \\[0.5em]
        \textbf{情感（$t-1$）} $\xrightarrow{-}$
        \textbf{讨论活动（$t$）}
        \vspace{0.5cm}
    }}
    \caption{研究二情感传导路径假设}
    \label{fig:path-model-2}
\end{figure}

\textit{注：箭头上的"$-$"号表示"情感极性下降"（即负面情感增加）对目标变量的正向影响，
或代码活动增加对讨论活动的抑制作用；
以情感极性的下降方向解读，方向符号与假设方向一致。}

\section{数据与变量}

\subsection{数据来源}

\begin{table}[htbp]
    \centering
    \caption{研究二数据来源}
    \label{tab:data-source-2}
    \begin{tabularx}{\textwidth}{lXl}
        \toprule
        数据类型 & 来源与说明 & 粒度 \\
        \midrule
        Issue 评论文本
            & GitHub API / GH Archive，Bitcoin 项目 Issue 区全量评论
            & 评论级 \\
        情感评分
            & \todo{说明所用工具：VADER / FinBERT / 其他}
            & 评论级→周级聚合 \\
        代码活动
            & GH Archive，CCE + PRE 事件总量
            & 周级 \\
        讨论活动
            & GH Archive，IE 事件总量
            & 周级 \\
        \bottomrule
    \end{tabularx}
\end{table}

\subsection{情感测量方法}

\todo{说明情感分析工具的选择依据、技术文本的处理挑战、
情感极性聚合方式（周均值或中位数）}

\subsection{关键变量}

\begin{table}[htbp]
    \centering
    \caption{研究二关键变量定义}
    \label{tab:variables-2}
    \begin{tabularx}{\textwidth}{llX}
        \toprule
        变量 & 类型 & 定义 \\
        \midrule
        $\text{Sentiment}_t$
            & 内生变量
            & 第 $t$ 周 Issue 评论情感极性均值，区间为 $[-1, 1]$，
              负值表示负面情感主导 \\
        $\text{CodeActivity}_t$
            & 内生变量
            & 第 $t$ 周 CCE + PRE 事件总量（对数） \\
        $\text{DiscussActivity}_t$
            & 内生变量
            & 第 $t$ 周 IE 事件总量（对数） \\
        \bottomrule
    \end{tabularx}
\end{table}

\section{研究方法：向量自回归模型（VAR）}

\subsection{模型选择的理由}

情感、代码活动和讨论活动三个变量之间存在潜在的双向影响关系，
不存在明确的先验主从结构，因此单方程回归不适合捕捉这种多元动态反馈。
向量自回归（Vector Autoregression, VAR）模型将所有变量均视为内生变量，
适合分析多变量时间序列系统中的动态因果关系。

\subsection{分析步骤}

\begin{enumerate}
    \item \textbf{平稳性检验：}对三个时间序列分别进行 ADF 单位根检验，
          若存在单位根则进行差分处理；
    \item \textbf{最优滞后阶数选择：}基于 AIC/BIC 准则确定最优滞后期数 $p$；
    \item \textbf{VAR($p$) 模型估计：}
          \begin{equation}
          \label{eq:var-model}
          \mathbf{y}_t = \mathbf{c} + \sum_{k=1}^{p} \mathbf{A}_k \mathbf{y}_{t-k} + \boldsymbol{\varepsilon}_t
          \end{equation}
          其中 $\mathbf{y}_t = [\text{Sentiment}_t,\,
          \text{CodeActivity}_t,\, \text{DiscussActivity}_t]'$，
          $\mathbf{A}_k$ 为系数矩阵，$\boldsymbol{\varepsilon}_t$ 为残差向量；
    \item \textbf{格兰杰因果检验（Granger Causality Test）：}
          检验各变量是否在统计意义上"格兰杰引起"其他变量；
    \item \textbf{脉冲响应函数（IRF）分析：}
          可视化某一变量受到一个标准差冲击后，其他变量随时间的动态响应路径；
    \item \textbf{方差分解（Forecast Error Variance Decomposition, FEVD）：}
          量化各变量对其他变量预测误差方差的贡献比例。
\end{enumerate}

\section{实证结果}

\subsection{平稳性检验}

\todo{ADF 检验结果表}

\subsection{最优滞后阶数}

\todo{AIC/BIC 信息准则结果，最终选择滞后 $p$ 期}

\subsection{格兰杰因果检验}

\todo{格兰杰因果检验结果表}

预期结果：
\begin{itemize}
    \item 情感 Granger-cause 代码活动（滞后1期）：显著 → H4
    \item 情感 Granger-cause 讨论活动（滞后1期）：显著 → H5
    \item 代码活动 Granger-cause 讨论活动（滞后2期）：显著 → H6
\end{itemize}

\subsection{脉冲响应分析}

\todo{IRF 图（3×3 矩阵或选取关键响应路径）及解读}

\subsection{方差分解}

\todo{FEVD 表及解读：情感对代码/讨论活动预测方差的贡献比例}

\section{稳健性检验}

\begin{itemize}
    \item 替换情感分析工具（VADER 与基于 BERT 的模型对比）
    \item 调整时间聚合粒度（月级 vs. 周级）
    \item 控制 Bitcoin 价格大幅波动周（异常值处理）
\end{itemize}

\todo{稳健性检验结果}

\section{小结}

本章基于情感即信息理论和目标手段替代理论，
采用 VAR 模型揭示了比特币社区情感对开发者活动的动态影响路径。
研究表明，负面情感作为项目问题的信号，
短期内激励开发者在代码和讨论两个维度上增加活动，
而代码活动的增加随后抑制了讨论活动，
揭示了"情感→代码/讨论→讨论"的传导链条。


%%====================================================================
\chapter{综合讨论}
%%====================================================================

\section{两项研究的共同发现}

本文两项研究从不同维度揭示了信息环境对开源开发者行为的塑造作用：
研究一聚焦外部信息输入（捐赠披露），研究二聚焦内部信息积累（社区情感）；
前者揭示信息激活期望与引发失望的双重效应，
后者揭示情感信号触发行动与活动间替代的动态机制。
两项研究的共同主题在于：\textbf{开发者的行为决策受到信息层面因素的深刻影响，
而非仅由实际物质激励驱动}。

\section{理论贡献}

\todo{完整阐述对激励理论、开源社区研究和图书情报学的理论贡献}

\section{实践启示}

\subsection{对开源基金会的启示}

\todo{针对 Brink 等机构的捐赠信息披露策略建议}

\subsection{对开源社区治理的启示}

\todo{情感监测与社区健康管理的启示}

\section{研究局限}

本研究存在以下局限：
\begin{itemize}
    \item 研究一中，Brink 定向披露事件数量有限，限制了参数估计的精度；
    \item 工具变量（Google 搜索指数）的排他性假设无法在统计层面完全验证；
    \item 情感分析工具在技术性文本（如代码讨论）上的准确性尚存局限；
    \item 观测期（2022—2023年）有限，结论向其他时期和项目的外推需谨慎；
    \item 未能区分情感的发出者（核心开发者 vs. 外围参与者）对结果的差异化影响。
\end{itemize}


%%====================================================================
\chapter{结论}
%%====================================================================

\section{主要研究结论}

本文以比特币开源社区为实证场景，围绕捐赠信息披露与社区讨论情感两个维度，
系统考察了信息因素对开发者活动的影响，得出以下主要结论：

\textbf{研究一结论：}
\begin{enumerate}
    \item 非定向捐赠披露通过提升开发者对奖励价值和可得性的期望，
          正向影响代码贡献与审查贡献（H1、H2 成立）；
    \item 定向捐赠披露通过强化"绩效—奖励"因果链的工具性感知，
          亦正向影响开发者贡献（H2 成立）；
    \item 定向披露公布后，未入选开发者经历期望失验，
          对非定向披露的响应强度显著减弱（H3 成立）；
    \item 上述效应独立于实际捐款金额的直接激励效应，
          系纯粹的信息披露效应。
\end{enumerate}

\textbf{研究二结论：}
\begin{enumerate}
    \item 社区情感极性下降（负面情感增加）在滞后一期正向影响代码活动与讨论活动（H4、H5 成立）；
    \item 代码活动增加在滞后两期负向影响讨论活动（H6 成立）；
    \item 存在"情感→代码/讨论→讨论"的完整影响路径，
          分别以情感即信息理论和目标手段替代理论加以解释。
\end{enumerate}

\section{研究展望}

未来研究可从以下方向延伸：
\begin{enumerate}
    \item 将研究场景拓展至以太坊（Ethereum）、Monero 等其他加密货币开源社区，
          开展跨社区比较研究；
    \item 结合开发者访谈或问卷调查，从主观认知层面直接验证心理机制；
    \item 细化情感类型分析，区分愤怒、沮丧、焦虑等负面情感的差异化影响；
    \item 将两项研究整合为统一分析框架，考察捐赠信息披露与社区情感的交互效应；
    \item 探索机器学习方法在更大规模开源社区中进行情感分析和行为预测的可行性。
\end{enumerate}

%---------------------------------------------------------------------
%   参考文献
%---------------------------------------------------------------------

\printbibliography

%---------------------------------------------------------------------
%   致谢
%---------------------------------------------------------------------

\begin{acknowledgement}
感谢导师的悉心指导与耐心支持。
感谢图书情报学院各位老师在研究方法和学术写作方面的宝贵建议。
感谢 Bitcoin 开源社区的全体开发者，正是他们的开放协作精神使本研究成为可能。
感谢 Bitcoin Brink 基金会公开其运营数据，为本研究提供了独特的研究场景。
感谢同门师兄弟姐妹在讨论中给予的灵感和鼓励。
\end{acknowledgement}

%---------------------------------------------------------------------
%   附录
%---------------------------------------------------------------------

\appendix

\chapter{Brink 捐赠信息披露事件列表}

\todo{整理2022—2023年期间所有非定向披露和定向披露事件的完整列表}

\chapter{研究一：2SLS 第一阶段回归结果}

\todo{插入第一阶段回归结果，说明工具变量相关性（F 统计量 > 10）}

\chapter{研究二：VAR 模型稳健性检验}

\todo{替换情感工具、调整时间粒度的稳健性检验结果}

\chapter{情感分析方法说明}

\todo{详细说明情感分析工具的技术细节及在比特币 Issue 文本上的验证}

\end{document}
